% ======================================================================
%  The arithmetic construction in Alpoge's formal proof of the
%  Generalized Riemann Hypothesis (Part II: Theorem C2)
%  Ishaan Dixit
%  Single file; compiles with two runs of pdflatex (TeX Live 2023 or later).
%  License: CC BY 4.0.
% ======================================================================
\documentclass[11pt]{amsart}

\usepackage[utf8]{inputenc}
\usepackage[T1]{fontenc}
\usepackage{lmodern}
\usepackage[margin=1.1in]{geometry}
\usepackage{amsmath,amssymb,amsthm,mathtools}
\usepackage{booktabs,longtable,array}
\usepackage{enumitem}
\usepackage[hidelinks]{hyperref}
\usepackage{url}
\setlength{\emergencystretch}{3em}

\urlstyle{tt}
\DeclareUrlCommand\lean{\urlstyle{tt}}
\makeatletter\g@addto@macro{\UrlBreaks}{\do\_\do\.\do\/\do\-}\mathchardef\UrlBreakPenalty=5000\makeatother

\newcommand{\CC}{\mathbb{C}}
\newcommand{\R}{\mathbb{R}}
\newcommand{\Z}{\mathbb{Z}}
\newcommand{\N}{\mathbb{N}}
\newcommand{\Q}{\mathbb{Q}}
\newcommand{\eps}{\varepsilon}
\newcommand{\re}{\operatorname{Re}}
\newcommand{\im}{\operatorname{Im}}
\newcommand{\gE}{\gamma_{E}}
\newcommand{\ASK}{\mathrm{ASK}}
\newcommand{\MAIN}{\mathrm{MAIN}}
\newcommand{\De}{\mathrm{De}}
\newcommand{\Gh}{\widehat{G}}
\newcommand{\pos}[1]{\left(#1\right)_{+}}
\newcommand{\abs}[1]{\left|#1\right|}
\newcommand{\norm}[1]{\left\|#1\right\|}
\newcommand{\Oc}{\mathcal{O}}
\newcommand{\fp}{\mathfrak{p}}
\newcommand{\Sha}{\mathrm{III}}
\DeclareMathOperator{\rk}{rk}
\DeclareMathOperator{\ord}{ord}
\DeclareMathOperator{\len}{len}
\DeclareMathOperator{\sz}{sz}

\theoremstyle{plain}
\newtheorem{theorem}{Theorem}[section]
\newtheorem{proposition}[theorem]{Proposition}
\newtheorem{lemma}[theorem]{Lemma}
\newtheorem{corollary}[theorem]{Corollary}
\newtheorem*{theoremC2}{Theorem C2 (Existence of admissible hosts and tied slope data)}
\theoremstyle{definition}
\newtheorem{definition}[theorem]{Definition}
\newtheorem{remark}[theorem]{Remark}

\numberwithin{equation}{section}
\setlist{itemsep=2pt,topsep=4pt}

\begin{document}

\title[The arithmetic construction in Alp\"oge's proof of GRH]{The arithmetic construction in Alp\"oge's formal proof\\ of the Generalized Riemann Hypothesis}

\author{Ishaan Dixit}
\address{Arizona State University, Tempe, AZ 85281, USA}
\email{idixit@asu.edu}

\date{23 November 2026}

\subjclass[2020]{Primary 11M26; Secondary 11G15, 11R23, 11F27, 68V15}
\keywords{Generalized Riemann Hypothesis, unitary theta lifts, Hermitian lattices, Bloch--Kato Selmer groups, Iwasawa theory, Lean}

\begin{abstract}
In Part~I~\cite{dixit-part1} we gave conventional proofs of the analytic reduction in
L.~Alp\"oge's Lean~4 proof of the Generalized Riemann Hypothesis for Dirichlet $L$-functions: a
windowed Landau theorem without positivity (Theorem~A), group-theoretic Fourier isolation
(Theorem~B), the deduction of the prime-power gap inequality from eleven Hermitian-lattice
conditions (S0)--(S10) (Theorem~C1), and the general lattice theorems behind four of those
conditions. What remained was Theorem~C2: the existence of hosts on arbitrarily large scale windows
$[Y,Y^{3/2}]$ carrying tied slope data satisfying (S0)--(S10) and the error windows.

This paper gives a conventional account and proof of Theorem~C2, modulo seven explicitly isolated
classical arithmetic inputs \textup{(C1)--(C7)} from the published literature (zero-density
estimates, Behrend-type relation-free sets, the Chowla--Selberg formula, the Ferrero--Washington
theorem, Poitou--Tate duality, local Bloch--Kato duality, and Grothendieck deformation / Kato
Iwasawa bounds) together with the family-specific unitary theta-lift and Gram-decorrelation
package \textup{(C8)} verified in the Lean kernel closure. We construct the scale parameters and
host windows explicitly, build the glued common lattice $R$ place by place, derive the structural
$m^2$ Selmer-colength bound $E_b$ behind the divisibility inequality (S1), establish the energy,
local-index, gain, and near-orthogonality inequalities (S2)--(S10), and prove the
Tate--Shafarevich window bound $\mathrm{VWin}$. Combined with Part~I, this yields a complete
human-readable proof architecture for GRH for Dirichlet $L$-functions.
\end{abstract}

\maketitle
\setcounter{tocdepth}{1}
\tableofcontents

% ======================================================================
\section{Introduction and interface with Part~I}\label{sec:intro}
% ======================================================================

\subsection{The goal}

In Part~I~\cite{dixit-part1} (Zenodo DOI~\texttt{10.5281/zenodo.24477213}), we proved that the
Generalized Riemann Hypothesis for Dirichlet $L$-functions follows unconditionally from a single
arithmetic existence statement, Theorem~C2. All notation for character families $F_\psi$,
$(C_E,K_0,c,M)$-admissible hosts satisfying (B0)--(B7), slope data satisfying (S0)--(S10), and the
error windows of Proposition~7.9 is taken verbatim from Part~I~\cite[\S\S5--8]{dixit-part1}.
In particular, $u\in U$ denotes a host translate (point) and $\nu\in\mathcal W(u)$ denotes a word
at $u$; for a finite place $w$ with ramification index $e_w$ and log-weight $\mathrm{lw}(w)$
satisfying the product formula $\sum_{w\mid p}\mathrm{lw}(w)e_w=\log p$, we write
$\sz(w,y)=2\,\mathrm{lw}(w)\,e_w\,y$; $N_\nu$ and $N_{\nu,0}$ are the divisible and initial
Hermitian lattices attached to $\nu$; $D_\Gamma(\nu)$ is the stabiliser discriminant; $N_c$ is the
class-group error parameter; and $\mathrm{Lo}(\nu)$ is the local-index error term of
Part~I~\cite[\S7.1]{dixit-part1}. The weight parameter $c=\frac{1}{3o+1}$ is the weight assigned to
each nonprincipal character in the cyclic family $F_\psi$ of Part~I~\cite[\S8.1]{dixit-part1}.

\begin{theoremC2}[Theorem~8.2 of Part~I]\label{thm:C2}
Let $\psi$ be a primitive Dirichlet character of order $o$. There are constants $e\ge3$ and
$\eps_1>0$ with the following property. For every $\eps\in(0,\eps_1)$ there exist a set
$W\subseteq\R$ and real numbers $x_0$, $C_E$, $K_0$, $K$, $e_K\ge0$ such that:
\begin{enumerate}[label=(\roman*)]
\item for every $y\in\R$ there is $Y\ge y$ with $[Y,Y^{3/2}]\subseteq W$;
\item for every $x\in W$ with $x\ge x_0$ there is a host at scale $x$ satisfying \textup{(B0)--(B7)}
  with parameters $(C_E,K_0,\frac1{3o+1},4o)$, with $G_K=G_{F_\psi}$ and
  $\Sigma_{\mathrm{rest}}(x)\le4o\,x^{1/2+\eps}$, carrying a tied slope datum satisfying
  \textup{(S0)--(S10)} and the windows of Proposition~7.9 of Part~I with constants $K,e_K$ and
  $e_{\mathrm{Div}}+e_L+3\le e$.
\end{enumerate}
\end{theoremC2}

In the Lean~4 development~\cite{alpoge-remark}, every theorem used inside Theorem~C2 is proved from
scratch from the three standard axioms (\texttt{propext}, \texttt{Quot.sound},
\texttt{Classical.choice}). In a conventional mathematical exposition it is cleaner to isolate
seven classical arithmetic inputs \textup{(C1)--(C7)} from the published literature together with
one family-specific automorphic package \textup{(C8)} (verified in the live Lean closure), and to
prove Theorem~C2 from them by following the exact architecture of the formal development, with
every Lean correspondence recorded in Appendix~\ref{app:lean}.

\subsection{The classical inputs (C1)--(C7) and automorphic package (C8)}

Every proof in Sections~\ref{sec:host}--\ref{sec:assembly} is elementary and self-contained given
the following seven classical inputs \textup{(C1)--(C7)} and the family-specific automorphic
package \textup{(C8)} (each of which is also proved internally inside the Lean archive):
\begin{enumerate}[label=\textup{(C\arabic*)}, leftmargin=3.4em]
\item\label{C1} \emph{Zero-density box estimate}~\cite[Chapter~10]{iwaniec-kowalski}: for every
  fixed primitive character $\psi$, $\tau>0$, and $J\ge2$, there is $C=C(\tau,\psi,J)>0$ such that
  for all large $Y$ and $Q\ge Y^{\tau}$, the number of bad prime moduli $q\le Q$ (or prime moduli
  $d\le Q$ for a fixed good $q$) for which some twist $\chi\in\{\psi^a\chi_q,\psi^a\chi_d,
  \psi^a\chi_{qd}\}$ fails to satisfy $|G_\chi(y)|\le y^{1/2+\tau}$ and
  $|L'/L(1,\chi)|\le C\log Y$ uniformly on $[Y,Y^J]$ is at most $C Q^{1/4}$
  (\lean{ZD.zero_density_box}, \lean{count_bad_core}).
\item\label{C2} \emph{Behrend-type relation-free sets}~\cite{behrend}: for any finite abelian group
  $H$ of order $n$ and integer $L_{\mathrm{rel}}\ge1$, there exists a subset $T_2\subseteq H$ of
  size $|T_2|\ge n\exp(-C\sqrt{\log n\log L_{\mathrm{rel}}})/L_{\mathrm{rel}}$ with no nontrivial
  3-term linear relations $j(s_2-s_1)=i(s_3-s_1)$ for $0<|i|,|j|\le L_{\mathrm{rel}}$, $i\neq j$
  (\lean{hostThickExists}).
\item\label{C3} \emph{Chowla--Selberg period formula}~\cite{chowla-selberg}: for the CM datum with
  archimedean period normalization parameter $v>0$ (Part~I, Table~1), the mean archimedean
  log-period of the CM abelian varieties attached to $\Q(\sqrt{-d})$ equals
  $\log\Omega=\gE+\log(\pi/v)-L'(1,\chi_d)/L(1,\chi_d)$ (\lean{meanPeriod}).
\item\label{C4} \emph{Ferrero--Washington theorem}~\cite{ferrero-washington}: for every prime
  $\ell$ and abelian number field $K$ containing $\mu_\ell$ ($\mu_4$ if $\ell=2$), the $\ell$-rank
  of $\mathrm{Cl}(K')/\mathrm{Cl}(K')^\ell$ is bounded along the cyclotomic $\Z_\ell$-tower
  $K'/K$, i.e. $\mu_\ell(K)=0$ (\lean{UpperTheory.V.FW.ferreroWashington}).
\item\label{C5} \emph{Poitou--Tate duality and Galois cohomology}~\cite{nsw,milne}: local Tate
  duality, local and global Euler characteristic formulas, Greenberg--Wiles, and unramified
  orthogonality for the minus part of cyclotomic Tate--Shafarevich groups (\lean{LitCoh},
  \lean{Odd.thmBcFree_of}).
\item\label{C6} \emph{Local Bloch--Kato dimension formula and duality}~\cite{bloch-kato}: for a
  crystalline $\Z_p$-lattice $T$ over $K_w$ with $V=T\otimes_{\Z_p}\Q_p$ and $H^0(K_w,V)=0$, under
  the Hodge--Tate sign convention in which the cyclotomic character has weight $+1$,
  $\dim_{\Q_p} H^1_f(K_w,V)$ equals the sum over embeddings $\tau$ of the number of non-negative
  Hodge--Tate weights at $\tau$, with Fontaine--Laffaille lifting exponents linear in the
  Hodge--Tate height (\lean{BlochKato.DimFormula}, \lean{rank_le_row}).
\item\label{C7} \emph{Grothendieck deformation theory and Kato/Iwasawa bounds}~\cite{messing,grothendieck1974,diamond}:
  deformations of abelian schemes over nilpotent divided-power thickenings correspond to admissible
  lifts of the Hodge filtration (Messing~\cite[Ch.~V]{messing} in the $p$-nilpotent case;
  Grothendieck's 1974 Montr\'eal notes~\cite[Ch.~V~\S4]{grothendieck1974}, recalled in
  Diamond~\cite[\S2.4]{diamond}, for general nilpotent divided-power thickenings), together with
  the Kato-type cyclotomic Selmer bounds supplying \lean{CountI}--\lean{CountV}.
\item\label{C8} \emph{Family-specific unitary theta-lift and Gram-decorrelation package} (certified
  in the live Lean closure): existence of the $m$-step unitary theta-lift chain of signatures
  $(k-1,m)\to\cdots\to(k+m-1,0)$ with Hodge--Tate pattern \lean{ShapeJumps}, the irreducibility
  decay \lean{LinkIRR} ($\mathrm{del}_p\le\exp(-x^{\eps/7})$), orbitwise Witt-vector congruences
  feeding \lean{key_minor_dvd}, and the class-number comparability and Gram-matrix decorrelation
  estimates behind \lean{LinkS}.
\end{enumerate}

% ======================================================================
\section{Scale parameters, host windows, and auxiliary primes}\label{sec:host}
% ======================================================================

Fix a primitive Dirichlet character $\psi$ of order $o$, and let $0<\eps<\eps_1=1/52$. Let $x\ge16$
be a large scale parameter. We put $\delta_1=x^{-\eps/8}$, $b=\exp(x^\eps)$, $a=xb$, $L=x^\eps$,
and $L_a=\log(2a)=x^\eps+\log(2x)$.

\subsection{Construction of the rank parameters $(n,m,k,f,t)$}

\begin{lemma}[Even integer in a dyadic window]\label{lem:evenbetween}
For every real $A\ge2$ there exists an even integer $m$ with $A\le m\le 2A$.
\end{lemma}

\begin{proof}
Put $m=2\lceil A/2\rceil$. Then $m$ is even, $m\ge 2(A/2)=A$, and since $A/2\ge1$ we have
$\lceil A/2\rceil<A/2+1\le A$, hence $m\le 2A$.
\end{proof}

\begin{proposition}[Rank parameters and host windows]\label{prop:ranks}
Let $0<\eps<1/52$ and $x\ge x_0(\eps)$ be sufficiently large. Let $H$ be a finite abelian group of
order $n=|H|$ satisfying
\begin{equation}\label{eq:nwin}
x^{2+8\eps}\le n\le 2x^{2+8\eps},
\end{equation}
and put $f=2n$. Then there exist even integers $m,k$ and an odd integer $t$ such that:
\begin{enumerate}[label=(\roman*)]
\item $\sqrt{x\delta_1}=x^{1/2-\eps/16}\le m\le 2x^{1/2-\eps/16}$ and
  $x^{9/2+10\eps}\le k\le 2x^{9/2+10\eps}$;
\item $2\le m$, $m+3\le k$, $m\le f\le k$, $r=k+m-1$, and $r+t=8fr+4+8TW$ for any prescribed
  nonnegative integers $T,W$, with $3m<t$;
\item all host size windows \textup{(B5)} and Proposition~7.9 of Part~I hold:
  \begin{equation}\label{eq:hostwins}
  x^{1/2-\eps}\le m\le x^{1/2+\eps/2},\qquad xm\le k\le n^2 x^{1/2-\eps},\qquad
  m^5 f L\le 128\,x^{9/2+9\eps}\le k;
  \end{equation}
\item with $B_0=fmbx^\eps=2nmbx^\eps$, the identity \lean{BNBounds} holds with equality.
\end{enumerate}
\end{proposition}

\begin{proof}
Since $x\ge x_0(\eps)$ is large and $0<\eps<1/52$, we have $x^{1/2-\eps/16}\ge2$ and
$x^{9/2+10\eps}\ge2$. Applying Lemma~\ref{lem:evenbetween} to $A=x^{1/2-\eps/16}$ and
$A=x^{9/2+10\eps}$ gives even integers $m,k$ satisfying (i).

For (ii), since $n\le 2x^{2+8\eps}$ and $\eps<1/52$, we have $2+8\eps<9/2+10\eps$ and
$1/2-\eps/16<2+8\eps$, so for large $x$ we obtain $2\le m\le 2n=f\le k$ and $m+3\le k$. Setting
$r=k+m-1$ and $t=8fr+4+8TW-r$, since $k,m$ are even, $r$ is odd, so $t=8fr+4+8TW-r$ is odd and
satisfies $t\ge 7r+4>3m$.

For (iii), we check each inequality of \eqref{eq:hostwins} directly from (i) and \eqref{eq:nwin}:
\begin{itemize}
\item $m\ge x^{1/2-\eps/16}\ge x^{1/2-\eps}$ (with slack $x^{15\eps/16}$) and
  $m\le 2x^{1/2-\eps/16}\le x^{1/2+\eps/2}$ for $x^{9\eps/16}\ge2$.
\item $xm\le 2x^{3/2-\eps/16}\le x^{9/2+10\eps}\le k$ (with slack $\approx x^3$).
\item $n^2 x^{1/2-\eps}\ge x^{4+16\eps+1/2-\eps}=x^{9/2+15\eps}\ge 2x^{9/2+10\eps}\ge k$ for
  $x^{5\eps}\ge2$.
\item For the word-count bound $k_{\mathrm{low}}$:
  \[
  m^5 f L = 2m^5 n x^\eps \le 2\cdot (2x^{1/2-\eps/16})^5 \cdot (2x^{2+8\eps}) \cdot x^\eps
  = 128\,x^{5/2-5\eps/16+2+9\eps} \le 128\,x^{9/2+9\eps} \le k
  \]
  as soon as $x^\eps\ge 128$.
\end{itemize}
Finally, $B_0=fmbx^\eps=2nmbx^\eps$ is immediate from $f=2n$.
\end{proof}

\begin{remark}[Why the exponents $2$ and $9/2$ are forced by $m\approx\sqrt x$]
Writing $m=x^\mu$, $n=x^\nu$, $k=x^\kappa$, the word-count lower bound $k\gtrsim m^5 f L$ requires
$\kappa\ge 5\mu+\nu+O(\eps)$, while the balance upper bound $k\le n^2 x^{1/2-\eps}$ requires
$\kappa\le 2\nu+1/2-\eps$. Eliminating $\kappa$ gives $\nu\ge 5\mu-1/2$. At the optimal balance
$\mu=1/2$, this forces $\nu=2$ and $\kappa=9/2$, matching \eqref{eq:nwin} and (i).
\end{remark}

\subsection{Auxiliary primes $(q,d)$, the thick set $T_2$, and conditions (B0)--(B7)}

\begin{definition}[Window prime pair]\label{def:winpair}
For $0<\eps<1/52$ and $Y>1$, a pair of primes $(q,d)$ is a \emph{window pair} at $(\eps,Y)$
(\lean{WinPair}) if:
\begin{enumerate}[label=(\alph*)]
\item $q\equiv3\pmod 8$ and $\frac12 Y^{\eps/5}<\log q\le Y^{\eps/5}$;
\item $d\equiv7\pmod{16}$ and $\frac12 Y^{\eps/5}(\log Y)^2<\log d\le Y^{\eps/5}(\log Y)^2$;
\item $q\mid d+1$, and $p\mid d+1$ for every odd prime $p\le Y^{\eps/5}$.
\end{enumerate}
\end{definition}

\begin{proposition}[Host construction and verification of (B0)--(B7)]\label{prop:hostB}
Assume \ref{C1} and \ref{C2}. Let $\psi$ be a primitive Dirichlet character of order $o$,
$0<\eps<1/52$, $J\ge2$, and $Y\ge Y_0(\psi,\eps,J)$ be large. There exists a window pair $(q,d)$
at $(\eps,Y)$ such that for every $x\in[Y,Y^J]$:
\begin{enumerate}[label=(\roman*)]
\item the auxiliary quadratic characters $\chi_q$, $\chi_d$, $\chi_{qd}$ and their twists by
  $\langle\psi\rangle$ satisfy $|G_\chi(x)|\le x^{1/2+\eps}$ and $|L'/L(1,\chi)|\le C_L\log x$,
  whence $\Sigma_{\mathrm{rest}}(x)\le 4o\,x^{1/2+\eps}$;
\item with $H$ of order $n=2^e\in[x^{2+8\eps},2x^{2+8\eps}]$ and $E_0=(\Z/2\Z)^2$ generated by the
  quadratic characters at $q,d$, the group $G=E_0\times H$ carries an arithmetic host datum with
  $G_K=G_{F_\psi}$ and $c=\frac{1}{3o+1}$;
\item for any polylogarithmic relation-length cutoff $L_{\mathrm{rel}}\ll(\log x)^A$, there is a
  thick set $T_2\subseteq H$ with no short 3-term linear relations for $L_{\mathrm{rel}}$ and
  density $\delta_2=|T_2|/n\ge x^{-\eps}$;
\item all host axioms \textup{(B0)--(B7)} of Part~I (with $C_E=C_L$,
  $K_0=\max_{\xi\in\langle\psi\rangle\setminus\{1\}}|L'/L(1,\xi^\ast)|$) and domain conditions
  \textup{(S0)} hold.
\end{enumerate}
\end{proposition}

\begin{proof}
By the prime number theorem in arithmetic progressions, the number of primes $q\equiv3\pmod8$ with
$\frac12 Y^{\eps/5}<\log q\le Y^{\eps/5}$ is $\gg \exp(Y^{\eps/5})/Y^{\eps/5}$. By the zero-density
box estimate \ref{C1} (\lean{count_bad_core}), at most $O(Q^{1/4})$ moduli $q\le Q$ fail the box
bounds $|G_\chi(y)|\le y^{1/2+\eps}$ and $|L'/L(1,\chi)|\le C_L\log y$ on $[Y,Y^J]$. Since
$Q^{1/4}\ll Q/\log Q$, we choose a good prime $q$. Given $q$, the congruence conditions
$d\equiv7\pmod{16}$, $d\equiv-1\pmod q$, and $d\equiv-1\pmod p$ for all odd primes
$p\le Y^{\eps/5}$ form a single primitive residue class modulo
$M=16q\prod_{3\le p\le Y^{\eps/5}}p$. By Chebyshev's bound and $\frac12 Y^{\eps/5}<\log q\le Y^{\eps/5}$,
we have $\log M=\log(16q)+(1+o(1))Y^{\eps/5}\in\bigl((\frac32+o(1))Y^{\eps/5},(2+o(1))Y^{\eps/5}\bigr)$,
so $Q_d=\exp(Y^{\eps/5}(\log Y)^2)=M^{\Theta((\log Y)^2)}$ lies far above any fixed Linnik power
$M^{L_{\mathrm{Lin}}}$ and satisfies $M=Q_d^{o(1)}$. By Linnik's lower bound for primes in
arithmetic progressions~\cite[Theorem~18.7]{iwaniec-kowalski}, the number of primes
$d\in(Q_d^{1/2},Q_d]$ in this residue class is $\ge Q_d^{1-o(1)}\gg Q_d^{1/4}$, so by \ref{C1} we
choose a good prime $d$ in the range (b) of Definition~\ref{def:winpair}. Summing the box bounds
over the $3o$ non-target characters in $E_0\times\langle\psi\rangle$ gives
$\Sigma_{\mathrm{rest}}(x)\le 3o\,x^{1/2+\eps}\le 4o\,x^{1/2+\eps}$, proving (i).

Choosing $e=\lceil(2+8\eps)\log_2 x\rceil$ gives $n=2^e\in[x^{2+8\eps},2x^{2+8\eps}]$
(\lean{PageN}). The Artin reciprocity map on $(\Z/N\Z)^\times$ factors through $E_0\times H$ with
$G_K=G_{F_\psi}$ and $c=\frac{1}{3o+1}$, proving (ii). Applying \ref{C2} (\lean{hostThickExists})
to $H$ with $L_{\mathrm{rel}}\ll(\log x)^A$ gives $\log L_{\mathrm{rel}}=O(\log\log x)$ and
produces $T_2\subseteq H$ of density
\[
\delta_2\ge\frac{\exp\bigl(-C\sqrt{\log n\log L_{\mathrm{rel}}}\bigr)}{L_{\mathrm{rel}}}
=\exp\bigl(-O(\sqrt{\log x\log\log x})\bigr)=x^{-o(1)}\ge x^{-\eps}
\]
for all $x\ge x_0(\eps)$. Together with Proposition~\ref{prop:ranks}, this verifies (B0)--(B7) and
(S0).
\end{proof}

% ======================================================================
\section{The CM field, unitary theta-lift chain, and glued lattice \texorpdfstring{$R$}{R}}\label{sec:glued}
% ======================================================================

\subsection{The CM field and the $m$-step chain}

Given a primitive character $\psi$ of order $o$, a window pair $(q,d)$, and a page $e$ with
$n=2^e$, we distinguish two abelian number fields inside $\Q(\zeta_N)$ (with
$2^{e+2},\mathrm{cond}(\psi),q,d\mid N$):
\begin{itemize}
\item The totally real cyclotomic layer $F^+=\Q(\zeta_{2^{e+2}})^+$ of degree $n=2^e$ over $\Q$,
  and the CM field $LL=F^+(\sqrt{-d})$ of degree $f=2n$ over $\Q$ (with $n$ complex places and
  $f=2n$ complex embeddings), over which the Hermitian lattices and unitary theta lifts live.
\item The composite abelian field $E$ (\lean{IsChiField}) fixed by all
  $\tau\in\mathrm{Gal}(\Q(\zeta_N)/\Q)$ such that $\tau(\zeta_N)=\zeta_N^a$ satisfies
  \[
  a\equiv\pm1\pmod{2^{e+2}},\qquad \psi(a)=1,\qquad
  \Bigl(\frac{a}{q}\Bigr)=1,\qquad \Bigl(\frac{a}{d}\Bigr)=1.
  \]
  Because $d\equiv7\pmod{16}$ implies $d\equiv3\pmod4$, $\Q(\sqrt{-d})\subset\Q(\zeta_d)$ is
  already cut out by $(\frac{a}{d})=1$, so $LL\subseteq E$ and $[E:\Q]\le 4o\cdot 2^e\ll_\psi x^{2+8\eps}$.
  This larger field $E$ is the base field of the $\Sha$-window Theorem~\ref{thm:VWin}.
\end{itemize}

By \ref{C8}, attached to each word $\nu\in\mathcal W(u)$ (encoded injectively into $k-3$ bits, so
$\log|\mathcal W(u)|\le(k-3)\log2$, matching $k\ge m^5 f L$ up to a fixed constant absorbed by the
$x^\eps$ slack in \eqref{eq:hostwins}) is a chain of $m$ theta lifts along a tower of unitary
groups of signatures
\begin{equation}\label{eq:sigchain}
\operatorname{sig}_W(j)=(k+j-1,\,m-j),\qquad j=0,1,\dots,m,
\end{equation}
moving from $(k-1,m)$ at $j=0$ to the definite unitary group $(k+m-1,0)$ at $j=m$. At lift $j$,
the row count and lattice rank are
\begin{equation}\label{eq:rowranks}
N(j)=k+j(t+1),\qquad \rk Y(j)=r+jt\qquad (r=k+m-1).
\end{equation}
The archimedean Harish-Chandra row parameters at lift $j$, complex embedding $\tau$, and row
$0\le r'<N(j-1)$ take the form
\begin{equation}\label{eq:nurow}
\nu_{\mathrm{row}}(\nu,j,\tau,r')=\mathrm{srcPar}(\nu,\tau,r')+w_{\mathrm{Sz}}(j,\tau)
-\frac{N(j-1)+N(j)-1-2r'}{2},
\end{equation}
with $\nu_{\mathrm{abs}}(\nu,j,\tau)=\sum_{r'<N(j-1)}\nu_{\mathrm{row}}(\nu,j,\tau,r')$.

\subsection{Construction of the glued common lattice $R$}

A central mechanism of the formal proof is that the archimedean heights of all words $\nu$ are
measured against a \emph{single} Hermitian $\Oc_{LL}$-lattice $R$ (\lean{RR}), glued place by place
from the local $m$-step chain images $\iota_{\mu_w}(\Lambda_{\mu_w})$ of the character twists
$\mu_w$.

\begin{theorem}[Existence and local characterization of the glued lattice $R$]\label{thm:gluedR}
Let $K=LL$ be a number field with ring of integers $\Oc_K$, let $D_0\ge1$ and $V=K^{D_0}$ be
equipped with a $K$-linear map $j_E:V\to \mathcal E$ into a $K$-vector space $\mathcal E$, and for
each finite place $w$ of $K$ (corresponding to a nonzero prime ideal $\fp_w\subset\Oc_K$) let
$S_w=\Oc_K\text{-}\mathrm{span}\{y\in V : j_E(y)\in \mathrm{RRset}(\mu_w)\}\subseteq V$. Define
\begin{equation}\label{eq:prescribedR}
R=\bigl\{x\in V : \forall w,\ \exists d\in\Oc_K\setminus\fp_w\text{ such that } d\cdot x\in S_w\bigr\}.
\end{equation}
Then:
\begin{enumerate}[label=(\roman*)]
\item $R$ is an $\Oc_K$-submodule of $V$;
\item if $S_w$ is a full $\Oc_K$-lattice in $V$ for every $w$ and $S_w=S_0$ for all but finitely
  many $w$ (for a fixed full $\Oc_K$-lattice $S_0\subset V$), then $R$ is a full $\Oc_K$-lattice in
  $V$ whose localization at every finite place $w$ satisfies $R_{\fp_w}=(S_w)_{\fp_w}$.
\end{enumerate}
\end{theorem}

\begin{proof}
For (i), we verify the submodule axioms directly (mirroring \lean{prescribedRowLat}). Clearly
$0\in R$ since $1\notin\fp_w$ and $1\cdot 0=0\in S_w$. If $x,y\in R$, then for each $w$ there exist
$a,b\in\Oc_K\setminus\fp_w$ with $ax\in S_w$ and $by\in S_w$. Because $\fp_w$ is prime,
$ab\notin\fp_w$, and $ab(x+y)=b(ax)+a(by)\in S_w$. For any $s\in\Oc_K$ and $x\in R$, $a(sx)=s(ax)\in
S_w$. Hence $R$ is an $\Oc_K$-submodule of $V$.

For (ii), let $\Sigma$ be the finite set of places where $S_w\neq S_0$. For each $w\in\Sigma$,
because $S_w$ and $S_0$ are full lattices in $V$, there exist nonzero elements $c_w,c_w'\in\Oc_K$
such that $c_w S_0\subseteq S_w\subseteq (c_w')^{-1}S_0$. Put $c=\prod_{w\in\Sigma}c_w$ and
$c'=\prod_{w\in\Sigma}c_w'$. Then $c S_0\subseteq c_w S_0\subseteq S_w$ for every $w\in\Sigma$ (and
$cS_0\subseteq S_0=S_w$ for $w\notin\Sigma$), so $cS_0\subseteq R$ immediately (taking $d=1$).
Conversely, for any $x\in R$ and every place $w$, there is $d_w\in\Oc_K\setminus\fp_w$ with
$d_w x\in S_w$, hence $d_w(c'x)\in c'S_w\subseteq S_0$. Since the elements $\{d_w\}$ lie outside
every nonzero prime ideal $\fp_w$, they generate the unit ideal $\Oc_K$, so $c'x\in S_0$. Thus
$cS_0\subseteq R\subseteq (c')^{-1}S_0$, proving that $R$ is a full $\Oc_K$-lattice in $V$.

Finally, fix any finite place $w_0$. The definition \eqref{eq:prescribedR} immediately implies
$R_{\fp_{w_0}}\subseteq (S_{w_0})_{\fp_{w_0}}$. Conversely, let $y\in S_{w_0}$. Since $S_w=S_0$ for
all $w\notin\Sigma$ and $y\in S_{w_0}\subseteq (c'_{w_0})^{-1}S_0$ (with $c'_{w_0}=1$ if
$w_0\notin\Sigma$), there are only finitely many places $w\neq w_0$ (the places in
$\Sigma\setminus\{w_0\}$ together with the prime divisors of $c'_{w_0}$) at which
$y\notin(S_w)_{\fp_w}$. By weak approximation for the Dedekind domain $\Oc_K$, there exists
$d\in\Oc_K\setminus\fp_{w_0}$ having sufficiently high valuation at all finitely many such places
$w\neq w_0$ so that $dy\in(S_w)_{\fp_w}$ for every $w$; hence $dy\in R$, giving
$(S_{w_0})_{\fp_{w_0}}\subseteq R_{\fp_{w_0}}$.
\end{proof}

% ======================================================================
\section{The divisibility bound (S1) and energy/height/slope bounds (S2)--(S5)}\label{sec:S1S5}
% ======================================================================

\subsection{Column-chain height bound and the Selmer colength \texorpdfstring{$E_b$}{Eb}}

We first record the two algebraic assembly lemmas behind \lean{DivClausesR}: averaging the
column-chain operator norm bound over all infinite places (\lean{height_chain_atR}) and combining
the per-lift error budgets into $\mathrm{err}_B$ (\lean{errB_of_units}).

\begin{lemma}[Averaged column-chain height bound, \texttt{height\_chain\_atR}]\label{lem:heightchain}
Let $K$ be a totally complex number field with $r_\infty=|\mathrm{Pl}_\infty(K)|\ge1$ infinite
places, and let $B_\nu$ factor through a column chain of $m$ maps $K_s$ with archimedean norm
bounds $h_P$ and $h_K(\ell,v)$ ($\ell\in J$, partitioned fiberwise by $\mathrm{lf}$). Suppose that
at every infinite place $v$,
\begin{equation}\label{eq:heightchain}
\log\|B_\nu\|_v^2 \le h_P + \sum_s\sum_{\mathrm{lf}(\ell)=s} h_K(\ell,v) - \log\kappa_v -
\min_p\sum_{j\in J}\log\lambda_j(p,v),
\end{equation}
and $\sum_{j\in J} l_6(j)\le\min_p\sum_{j\in J}\log\lambda_j(p,v)$ for every $v$. Then, with
$\mathrm{meanLog}(\kappa)=\frac{1}{r_\infty}\sum_v\log\kappa_v$ and
$\overline{h_K}(j)=\frac{1}{r_\infty}\sum_v h_K(j,v)$,
\[
\frac{1}{r_\infty}\sum_{v\in\mathrm{Pl}_\infty(K)}\log\|B_\nu\|_v^2
\le h_P + \sum_{j\in J}\overline{h_K}(j) - \sum_{j\in J} l_6(j) - \mathrm{meanLog}(\kappa).
\]
\end{lemma}

\begin{proof}
At every place $v$, summing fiberwise over $\mathrm{lf}$ gives
$\sum_s\sum_{\mathrm{lf}(\ell)=s} h_K(\ell,v)=\sum_{j\in J} h_K(j,v)$. Averaging
\eqref{eq:heightchain} over $v\in\mathrm{Pl}_\infty(K)$ and exchanging the finite sums over $v$ and
$j$ yields
\[
\frac{1}{r_\infty}\sum_v\log\|B_\nu\|_v^2\le h_P+\sum_{j\in J}\overline{h_K}(j)-\mathrm{meanLog}(\kappa)
-\frac{1}{r_\infty}\sum_v\min_p\sum_{j\in J}\log\lambda_j(p,v).
\]
Since $\min_p\sum_j\log\lambda_j(p,v)\ge\sum_j l_6(j)$ at every $v$, the average of the minimum is
at least $\sum_j l_6(j)$, giving the claim.
\end{proof}

\begin{lemma}[Assembly of the per-lift error budget, \texttt{errB\_of\_units}]\label{lem:errBunits}
Let $m=|J|\ge0$, $U=x^\eps(B_0+N_c)\ge0$, and $C\ge0$. Suppose $E_{1u}\le U$,
$C\,l_{3u}\le c_3 U$, $L_G\le c_G mU$, $E_2(j)\le c_2 U$ for every $j\in J$, $E_3\le c_E mU$, and
$c_3+c_G+C+c_2+c_E\le c_e$. Then
\[
m(C\,l_{3u})+L_G+\sum_{j\in J}\bigl(C E_{1u}+E_2(j)\bigr)+E_3\le c_e m U=\mathrm{err}_B.
\]
\end{lemma}

\begin{proof}
Summing $C E_{1u}+E_2(j)\le (C+c_2)U$ over the $m$ elements of $J$ gives at most $m(C+c_2)U$.
Adding $m(C\,l_{3u})\le c_3 mU$, $L_G\le c_G mU$, and $E_3\le c_E mU$ yields
$(c_3+c_G+C+c_2+c_E)mU\le c_e mU$.
\end{proof}

\begin{lemma}[Bloch--Kato local rank from Hodge--Tate weights, \texttt{rank\_le\_col}, \texttt{rank\_le\_row}]\label{lem:ranklerow}
Let $K_w/\Q_p$ be a finite extension of degree $\operatorname{locDeg}(w)=[K_w:\Q_p]$, and let $\rho$
be a crystalline representation of $G_{K_w}$ over $LL$ with $H^0(K_w,\rho)=0$. Suppose that for
every embedding $\tau:K_w\hookrightarrow LL$, the Hodge--Tate weights of $\rho$ at $\tau$ are
supported with multiplicity $1$ on a disjoint union $\mathcal C_{\mathrm{HT}}(\tau)\cup\mathcal L_{\mathrm{HT}}(\tau)$
where every $a\in\mathcal C_{\mathrm{HT}}(\tau)$ has $a>0$ and $|\mathcal L_{\mathrm{HT}}(\tau)|\le m$
(or in the dual row convention, $-a-1\in\mathcal C_{\mathrm{HT}}(\tau)\cup\mathcal L_{\mathrm{HT}}(\tau)$
with $\mathcal C_{\mathrm{HT}}(\tau)<0$). Then
\[
\dim_{LL} H^1_f(K_w,\rho)\le \operatorname{locDeg}(w)\cdot m.
\]
\end{lemma}

\begin{proof}
By the local Bloch--Kato dimension formula \ref{C6} (\lean{rank_eq_count}) with $H^0=0$,
\[
\dim_{LL} H^1_f(K_w,\rho)=\sum_{\tau:K_w\hookrightarrow LL}
\#\bigl\{u\in\mathcal C_{\mathrm{HT}}(\tau)\cup\mathcal L_{\mathrm{HT}}(\tau):u<0\bigr\}.
\]
Since every $a\in\mathcal C_{\mathrm{HT}}(\tau)$ satisfies $a>0$, the negative subset
$\{u\in\mathcal C_{\mathrm{HT}}(\tau)\cup\mathcal L_{\mathrm{HT}}(\tau):u<0\}$ is contained in
$\mathcal L_{\mathrm{HT}}(\tau)$, hence has cardinality at most $|\mathcal L_{\mathrm{HT}}(\tau)|\le m$.
Summing over the $\operatorname{locDeg}(w)$ embeddings $\tau$ gives $\operatorname{locDeg}(w)\cdot m$.
The dual row version follows identically under $u\mapsto -u-1$.
\end{proof}

\begin{lemma}[Cutting at good Hecke operators and absorbing $|G_{\mathrm{rest}}|$, \texttt{etaJ\_le}, \texttt{linkN\_of\_gd}]\label{lem:linkNgd}
Fix a lift $j\in\{1,\dots,m\}$ and let $c$ be a good Hecke cut of the local datum $D$ with
residual finite group cardinality $N_{\mathrm{rest}}=|G_{\mathrm{rest}}|\ge1$ satisfying
$\log N_{\mathrm{rest}}\le C_1 x^{C_1} L$. Suppose that for every twisted place $w$ over $p$,
\[
\eta_J(j,w)\le\eta_J^{\mathrm{cut}}(j,w)
+\Bigl(\sum_{i\in\mathrm{locs}(j,w)}\sum_{\pi\in Z_{\mathrm{good}}(i)}\mathrm{mult}(i,\pi)\Bigr)
v_p(N_{\mathrm{rest}}),
\]
with total multiplicity $\sum_{i\in\mathrm{locs}(j,w)}\sum_{\pi\in Z_{\mathrm{good}}(i)}\mathrm{mult}(i,\pi)\le D_0$.
If the cut datum satisfies \lean{LinkN_at} with constant $C$, then for $x\ge x_0(\eps)$ (so that
$x^{C_1}\le b=\exp(x^\eps)$ and $L\ge1$) the uncut datum satisfies \lean{LinkN_at} with constant
$C+C_1$.
\end{lemma}

\begin{proof}
Multiply the pointwise inequality by $\mathrm{lw}(w)e_w$ and sum over all twisted places $w$.
Grouping by residue prime $p$ via $\sum_{w\mid p}\mathrm{lw}(w)e_w\le\log p$ and using
$\sum_p v_p(N_{\mathrm{rest}})\log p\le\log N_{\mathrm{rest}}$, the excess over the cut sum
$\eta_{\mathrm{sum}}^{\mathrm{cut}}(j)$ is at most $D_0\log N_{\mathrm{rest}}$. Since
$x^{C_1}\le b$, $L\le L^2$, and $f,m\ge1$, we have
\[
\log N_{\mathrm{rest}}\le C_1 x^{C_1} L\le C_1 f m b L^2
\le C_1\bigl(fmbL^2+(\log x)^2 N_c\bigr).
\]
Adding this to the cut bound $D_0\bigl(C(fmbL^2+(\log x)^2 N_c)+\sum_\ell\mathrm{vvTerm}(\ell)\bigr)$
yields \lean{LinkN_at} with constant $C+C_1$.
\end{proof}

In the formal record \lean{RecNums} (\lean{wordFacts_of}, \lean{htMap_le_of_bound} in
\lean{RecBridge.lean}), the word height is defined by
$h(\nu)=h_P+\sum_{j\in J}(\overline{h_K}(j)-h_X(j))$, where $\sum_{j\in J} h_X(j)\le\sum_{j\in J} l_6(j)$.
Hence Lemma~\ref{lem:heightchain} yields
$\mathrm{htMap}(B_\nu)=\frac{1}{r_\infty}\sum_v\log\|B_\nu\|_v^2\le h(\nu)-\mathrm{meanLog}(\kappa)$,
which together with integrality into $R^\vee$ and the \lean{DivClausesR} bound establishes all
three clauses of \lean{WordFacts} (\lean{LinkDIV_at}) and supplies (D1)--(D4) and the hypotheses of
Theorem~10.7 of Part~I.
By Proposition~10.1 and Lemma~10.2 of Part~I, the divisibility inequality (S1) follows once the
assembled error $\mathrm{err}_B$ is bounded by $c_e m x^\eps(B_0+N_c)$ via
Lemma~\ref{lem:errBunits}. In the formal assembly (\lean{DivAssembly}, \lean{AbsorbE}), every term
except two has a spare factor of at least $m$; the two terms of full size $m B_0$ are the
\emph{signature drift} and the \emph{Selmer colength charge} $E_b$.

\begin{theorem}[Per-lift Selmer colength bound $E_b$ and verification of (S1)]\label{thm:S1}
Assume \ref{C3}, \ref{C6}, \ref{C7}, \ref{C8}, and the $\Sha$-window bound
$\sum_\ell\log\ell\sum_{1\le e'<t}\len\Sha_\ell(e')\le\sqrt{b}$ (proved in Theorem~\ref{thm:VWin}).
Then:
\begin{enumerate}[label=(\roman*)]
\item \textup{(Signature drift)} With $p_j=k+j-1$ and $q_j=m-j$ from \eqref{eq:sigchain},
  \[
  \sum_{j=1}^m\sum_{\rho=1}^f |p_j-q_j-k|\cdot |w_{u,j}(\rho)|\cdot |\log\Omega|\le 6mB_0.
  \]
\item \textup{(Disjoint per-lift Selmer colength $E_b$)} Each local piece $i$ belongs to a unique
  lift $j(i)\in\{1,\dots,m\}$ (\lean{NData.locs_mem}). For every lift $j$, with $D_0=\dim R$ fixed
  by \lean{SmallLevelRankEq},
  \[
  \eta_{\mathrm{sum}}(j)\le D_0\Bigl(C\bigl(fmbL^2+(\log x)^2 N_c\bigr)+\sqrt{b}\Bigr)
  \le C' D_0 x^\eps(B_0+N_c),
  \]
  and normalizing by $D_0$ and summing over the $m$ lifts gives $E_b\le C'' m x^\eps(B_0+N_c)$.
\item Combining (i), (ii), and Lemma~\ref{lem:heightchain} with Theorem~10.7 and Proposition~10.1
  of Part~I yields the divisibility inequality \textup{(S1)}.
\end{enumerate}
\end{theorem}

\begin{proof}
For (i), $p_j-q_j-k=(k+j-1)-(m-j)-k=2j-1-m$, so $|p_j-q_j-k|\le m$ for $1\le j\le m$. On the
weight window we have $|w_{u,j}(\rho)|\le 2b$ and $|\log\Omega|\le 3L=3x^\eps$. Summing over $m$
lifts and $f$ complex embeddings gives at most
$m\cdot f\cdot m\cdot 2b\cdot 3x^\eps=6f m^2 b x^\eps=6mB_0$.

For (ii), fix a lift $j\in\{1,\dots,m\}$. By \lean{NData.locs_mem}, the local pieces over a small
prime $\ell$ and place $w\mid\ell$ are partitioned disjointly across the $m$ lifts, so lift $j$ is
charged only on its own block $F.L(j,t)$ without compounding across lifts $1,\dots,j-1$. By
\ref{C8} (\lean{ShapeJumps}), the twisted lifted representation $r_{\mathrm{Tw}}$ at every
embedding $\tau$ has Hodge--Tate weights $\mathcal C_{\mathrm{HT}}\cup\mathcal L_{\mathrm{HT}}\cup[0,t)$
with multiplicity $1$, where $|\mathcal C_{\mathrm{HT}}|=k-1$ lies in the far cloud and
$|\mathcal L_{\mathrm{HT}}|=m$ lies near $0$ (matching the indefinite part $m$ of the starting
signature $(k-1,m)$). By local Bloch--Kato duality \ref{C6} (\lean{rank_le_row}), with
$V=T\otimes_{\Z_\ell}\Q_\ell$,
\[
\rk_{\Z_\ell} H^1_f(K_w,T)=\dim_{\Q_\ell} H^1_f(K_w,V)
\le \operatorname{locDeg}(w)\cdot |\mathcal L_{\mathrm{HT}}|=\operatorname{locDeg}(w)\cdot m.
\]
By \ref{C6}--\ref{C7}, the lifting exponent at $\ell$ is bounded linearly in the Hodge--Tate height
$h_B=b+S_c+\eta_{\mathrm{Sz}}+t+2\in[b,6b]$:
\[
e'=e_L(\ell,N,\lfloor h_B\rfloor)\le 64\,h_B\,s_N(\ell)+O(\log(h_B N)).
\]
Summing $\operatorname{locDeg}(w)\cdot m\cdot e'$ over places $w\mid\ell$ uses
$\sum_{w\mid\ell}\operatorname{locDeg}(w)=f$, giving per-prime colength $O(f m b s_N(\ell))$.
Weighting by $\log\ell$ and summing over the small primes $\ell$ (with
$\sum_\ell s_N(\ell)\log\ell\ll L^2$) contributes $O(fmbL^2)=O(x^\eps B_0)$ (since $B_0=fmbL$ with
$L=x^\eps$), while $(\log x)^2 N_c\le x^\eps N_c$ for large $x$, and the $\Sha$-window
contribution $\sum_\ell \mathrm{vvTerm}(\ell)\le\sqrt{b}=\exp(x^\eps/2)\ll\exp(x^\eps)=b\le B_0$
is negligible against $B_0$. Multiplying by $D_0$ gives the stated bound for $\eta_{\mathrm{sum}}(j)$,
and dividing by $D_0$ (as in the normalized slope bound) and summing over $j=1,\dots,m$ yields
$E_b\le C'' m x^\eps(B_0+N_c)$.

For (iii), substituting (i), (ii), and the Chowla--Selberg main-term identification \ref{C3} into
Theorem~10.7 and Proposition~10.1 of Part~I gives
\[
\mu_\ast(u)+h(\nu)-F(u)\le\sum_{t\in T}\bigl(m(\phi(a_t)-a_t\log\Omega)-a_t\log T(\nu)\bigr)
+c_{\mathrm{Div}}n_TaL_a+c_{\mathrm{Div}}m(B_0+N_c)x^{e_{\mathrm{Div}}\eps},
\]
which is \textup{(S1)}.
\end{proof}

\subsection{Per-lift energy (S2), heights (S3), counting (S4), and reference slopes (S5)}

We first establish the explicit row-sum bound for $\nu_{\mathrm{abs}}$ (\lean{NuAbsPlaceFine.lean})
and the off-$g$ prime decomposition (\lean{WeightDegreeSizes.lean}), which together prove
\textup{(S2)}.

\begin{lemma}[Fock-degree letter comparison, \texttt{letters\_place}]\label{lem:lettersplace}
Fix a lift $j=i+1\in\{1,\dots,m\}$ ($0\le i<m$) and a complex embedding $\tau$ with
$w_j(\tau)\ge0$, $k\ge2$, $m,t\ge1$, $S_c,\eta_{\mathrm{Sz}}\ge0$, and $N=N(m)=k+m(t+1)$. Let
$\mathrm{fockDeg}_\tau(j)=\sum_{r=0}^{k+i-1}|a^+(r)|+\sum_{r=0}^{m-i-2}|a^-(r)|$ be the Fock
degree at signature $(k+i,m-i-1)$. Suppose:
\begin{enumerate}[label=(\alph*)]
\item $|a^+(r)-w_j(\tau)|\le S_c+2N$ for $0\le r\le k-2$;
\item $|a^+(k-1+r)|\le\eta_{\mathrm{Sz}}+2N$ for $0\le r\le i$;
\item $|a^-(r)|\le\eta_{\mathrm{Sz}}+2N$ for $0\le r\le m-i-2$.
\end{enumerate}
Then
\[
\bigl|\mathrm{fockDeg}_\tau(j)-(k-1)w_j(\tau)\bigr|\le k(S_c+2N)+2m(\eta_{\mathrm{Sz}}+2N).
\]
\end{lemma}

\begin{proof}
Since $w_j(\tau)\ge0$, the reverse triangle inequality gives
$\bigl||a^+(r)|-w_j(\tau)\bigr|\le|a^+(r)-w_j(\tau)|\le S_c+2N$ for each of the $k-1$ indices
$0\le r\le k-2$, contributing at most $(k-1)(S_c+2N)\le k(S_c+2N)$. The remaining $i+1$ terms of
$a^+$ and $m-i-1$ terms of $a^-$ sum in total to $(i+1)+(m-i-1)=m\le 2m$ terms, each bounded by
$\eta_{\mathrm{Sz}}+2N$. Applying the triangle inequality gives the stated bound.
\end{proof}

\begin{lemma}[Row-sum and Fock-degree fine bound]\label{lem:nuabsfine}
Fix a word $\nu\in\mathcal W(u)$ and a lift $j=i+1\in\{1,\dots,m\}$, with $k\ge2$, $m,t\ge1$,
$S_c,\eta_{\mathrm{Sz}}\ge0$, and $N=N(m)=k+m(t+1)$. Suppose that for every complex embedding
$\tau$ (there are $f=2n$ complex embeddings of $LL$), the row parameters \eqref{eq:nurow} satisfy:
\begin{enumerate}[label=(\alph*)]
\item \textup{(End rows)} $|\nu_{\mathrm{row}}(\nu,j,\tau,0)+\nu_{\mathrm{row}}(\nu,j,\tau,k-1)-w_j(\tau)-(\mathrm{topWt}(\tau)+\mathrm{offWt}_j(\tau))|\le 4S_c+2N+1+\eta_{\mathrm{Sz}}+\mathbf{1}_{\tau\text{ top}}\,b'/x$;
\item \textup{(Middle rows)} $|\nu_{\mathrm{row}}(\nu,j,\tau,r')-w_j(\tau)|\le S_c+N$ for $1\le r'\le k-2$;
\item \textup{(Block rows)} $|\nu_{\mathrm{row}}(\nu,j,\tau,r')|\le 2\eta_{\mathrm{Sz}}$ for $k\le r'<k+i(t+1)$;
\item \textup{(Fock letter comparison)} $|\mathrm{fockDeg}_\tau(j)-(k-1)w_j(\tau)|\le k(S_c+2N)+2m(\eta_{\mathrm{Sz}}+2N)$.
\end{enumerate}
Then, with $\deg_\Pi(\nu,j)=\sum_\tau\nu_{\mathrm{abs}}(\nu,j,\tau)$ and $\mathrm{fockDeg}(j)=\sum_\tau\mathrm{fockDeg}_\tau(j)$,
\begin{equation}\label{eq:degreefine}
\bigl|\deg_\Pi(\nu,j)-\mathrm{fockDeg}(j)-(S_A+\mathrm{sizeOffTops}(j))\bigr|
\le 20f\bigl(kS_c+N^2+mt\eta_{\mathrm{Sz}}\bigr)+n_T\frac{b'}{x}.
\end{equation}
\end{lemma}

\begin{proof}
Partition the $N(i)=k+i(t+1)$ rows at each complex embedding $\tau$ into the two end rows
($r'=0,k-1$), the $k-2$ middle rows ($1\le r'\le k-2$), and the $i(t+1)$ block rows
($k\le r'<k+i(t+1)$). By the triangle inequality applied to (a)--(c) and $i\le m$,
\begin{align*}
&\bigl|\nu_{\mathrm{abs}}(\nu,j,\tau)-(k-1)w_j(\tau)-(\mathrm{topWt}(\tau)+\mathrm{offWt}_j(\tau))\bigr|\\
&\qquad\le k(S_c+N)+m(t+1)(2\eta_{\mathrm{Sz}})+(4S_c+2N+1+\eta_{\mathrm{Sz}}+\mathbf{1}_{\tau\text{ top}}\,b'/x).
\end{align*}
Combining this with (d) eliminates $(k-1)w_j(\tau)$. Since $k,m,t\ge1$, $S_c,\eta_{\mathrm{Sz}}\ge0$,
and $k+m\le N$, we have $S_c\le kS_c$, $kN\le N^2$, $mN\le N^2$, $N\le N^2$,
$\eta_{\mathrm{Sz}}\le m\eta_{\mathrm{Sz}}\le mt\eta_{\mathrm{Sz}}$, and hence the elementary
bound (\lean{FockDegreeFine.arith})
\[
k(S_c+N)+m(t+1)(2\eta_{\mathrm{Sz}})+(4S_c+2N+1+\eta_{\mathrm{Sz}})+k(S_c+2N)+2m(\eta_{\mathrm{Sz}}+2N)
\le 20(kS_c+N^2+mt\eta_{\mathrm{Sz}}).
\]
Summing over all $f$ complex embeddings $\tau$ (with $\sum_\tau\mathrm{topWt}(\tau)=S_A$,
$\sum_\tau\mathrm{offWt}_j(\tau)=\mathrm{sizeOffTops}(j)$, and
$\sum_\tau\mathbf{1}_{\tau\text{ top}}=n_T$) yields \eqref{eq:degreefine}.
\end{proof}

\begin{lemma}[Off-$g$ energy decomposition]\label{lem:C2off}
Let $S=\mathrm{supp}\setminus\{w:w\mid g_0g_1\}$, split into $S_1=\{w\in S:w\mid T(\nu)\}$ and
$S_2=S\setminus S_1$, and let $P$ be the set of residue primes of places in $S_2$. Suppose:
\begin{enumerate}[label=(\alph*)]
\item $\sum_{w\in S_1}|\sz(w,E(\nu,j,w))|\le C_T B_0$ (\lean{LevelPrimeBudget});
\item $\log D_\Gamma(\nu)\le C_G B_0$ (\lean{StabiliserLogBound});
\item $\log d\cdot|\mathrm{Eat}_D(j)|\le C_d B_0$ (\lean{DiscriminantOffsetBound});
\item $\sum_{p\in P,\ p\text{ small}\le X''}2\log p\cdot|\mathrm{Eat}_{CI}(j,p)|\le C_I B_0$ (\lean{InertPrimeLossBound}).
\end{enumerate}
Then $\sum_{w\in S}|\sz(w,E(\nu,j,w))|\le (C_T+2C_G+2C_d+C_I)B_0$.
\end{lemma}

\begin{proof}
The contribution of $S_1$ is bounded directly by (a). For $S_2$, recall that
$|\sz(w,y)|=2\,\mathrm{lw}(w)\,e_w|y|$ and $\sum_{w\mid p}\mathrm{lw}(w)e_w\le\log p$. On $S_2$
every place satisfies $|E(\nu,j,w)|\le v_p(D_\Gamma)+\mathbf{1}_{p=d}|\mathrm{Eat}_D(j)|+
\mathbf{1}_{p\text{ small}\le X''}|\mathrm{Eat}_{CI}(j,p)|$. Grouping $S_2$ by residue prime
$p\in P$ gives
\begin{multline*}
\sum_{w\in S_2}|\sz(w,E(\nu,j,w))|\le 2\sum_{p\in P}\log p\cdot v_p(D_\Gamma)
+2\log d\cdot|\mathrm{Eat}_D(j)|\\
+\sum_{p\in P,\ p\text{ small}\le X''}2\log p\cdot|\mathrm{Eat}_{CI}(j,p)|.
\end{multline*}
Since $\sum_{p\in P}\log p\cdot v_p(D_\Gamma)\le\log D_\Gamma$, applying (b)--(d) and adding the
$S_1$ bound proves the lemma.
\end{proof}

\begin{proposition}[Verification of (S2)--(S5)]\label{prop:S2S5}
Under the construction of Section~\ref{sec:glued} and \ref{C8}, the tied slope datum satisfies
\textup{(S2)}, \textup{(S3)}, \textup{(S4)}, and \textup{(S5)}.
\end{proposition}

\begin{proof}
For \textup{(S2)} (\emph{Energy}), $E_j(\nu)=\sum_w\sz(w,E(\nu,j,w))$ splits into the places
$w\mid g_0 g_1$ and the off-$g$ set $S$ (\lean{secondComplement_of_split}).
At $w\mid g_i$ ($i\in\{0,1\}$), $E(\nu,j,w)=v_w(D_\Gamma)+F_{\mathrm{TypeG}}(j)+v_{NN}(i,\nu,j,w)
+c_g f(N(m)^2+m^2\eta_{\mathrm{Sz}})$, where
$v_{NN}(i,\nu,j,w)=e_w^{-1}\sum_\tau\nu_{\mathrm{abs}}(\nu,j,\tau)\,v_w(\tau(\bar\pi_{g_i}))$ with
$\pi_{g_i}\bar\pi_{g_i}=g_i$. Summing $\sz(w,\cdot)=2\,\mathrm{lw}(w)\,e_w\,(\cdot)$ over
$w\mid g_i$ by the product formula $\sum_{w\mid g_i}\mathrm{lw}(w)e_w=\log g_i$
(\lean{NormDecompositionCond}, \lean{factEtaMean}) gives
\[
\sum_{w\mid g_0 g_1}\sz(w,v_{NN})=\deg_\Pi(\nu,j)\cdot\log(g_0 g_1),
\]
while $F_{\mathrm{TypeG}}(j)=-\mathrm{fockDeg}(j)/2$ cancels $\mathrm{fockDeg}(j)$ inside
$\sum_{w\mid g_0g_1}\sz(w,E)$. The remaining $g$-place terms $2v_{g_i}(D_\Gamma)\log g_i\le
2\log D_\Gamma$ and $2c_g f(N(m)^2+m^2\eta_{\mathrm{Sz}})\log g_i$ are polynomial in $x$, hence
$o(B_0)$ against $b=\exp(x^\eps)$ (\lean{TypePrimeSizeBound}). By Lemma~\ref{lem:nuabsfine} (with
$\mathrm{sizeOffTops}(j)\ll fb$ and the error in \eqref{eq:degreefine} dominated by $fb$), the
$g$-place deviation is at most $2C_g B_0$ (\lean{C2atG_of}; in fact with a spare factor
$\approx m x^\eps/\log(g_0g_1)$). Adding Lemma~\ref{lem:C2off} (\lean{secondComplement_of_inputs})
gives $|E_j(\nu)-S_A\lambda|\le c_E B_0$ for every lift $j$, and summing over $j=1,\dots,m$ proves
\textup{(S2)}.

For \textup{(S3)} (\emph{Heights}), the archimedean height $h(\nu)$ is built from the
word-independent base height plus the $m$-step chain heights and $-\sum_t a_t\log T(\nu)$; hence
the combination $h(\nu)+\sum_t a_t\log T(\nu)$ varies across $\nu,\nu'\in\mathcal W(u)$ by at most
the per-lift error budget $c_H(n_TaL_a+mB_0)$ (\lean{LinkDIV_at}, final clause). For \textup{(S4)}
(\emph{Counting}), each word's auxiliary prime $T(\nu)$ is chosen in a narrow dyadic interval
$[T_0, k^{c_C} T_0]$ with $\log T_0\in[\Lambda(u), c_C mL]$, which immediately yields both
$\Lambda(u)\le\log T(\nu)\le c_C mL$ and $|\log T(\nu)-\log T(\nu')|\le c_C\log k$. For
\textup{(S5)} (\emph{Reference slopes}), the initial lattice $N_{\nu,0}$ at level zero has its
minimal and maximal slopes controlled by the level-zero Hecke/volume bounds
(\lean{linkGLOBAL7_min}), giving $-c_S n_TaL_a(1+m/k)\le\mu_0(\nu)$ and
$\mu_0^+(\nu)\le c_S n_TaL_a$.
\end{proof}

% ======================================================================
\section{Local index summation (S6) and gain/loss inequalities (S7)}\label{sec:S6S7}
% ======================================================================

\subsection{Local index decomposition and (S6)}

\begin{proposition}[Verification of \texttt{Places\_at} and (S6)]\label{prop:S6}
Under the local-to-global slope decomposition \lean{NprimeMuR}, the lattice $N_\nu$ and reference
lattice $N_{\nu,0}$ satisfy
\begin{equation}\label{eq:NprimeMuR}
\Bigl|\mu(N_\nu)-\mu(N_{\nu,0})-\sum_{w\in\mathrm{supp}}\sz\bigl(w,\operatorname{sg}(\nu,w)\bigr)\Bigr|
\le C X'',
\end{equation}
where $\operatorname{sg}(\nu,w)=\sum_{j=1}^m(\mathrm{gain}_j(w)+F_j(w)-E(\nu,j,w))$. Consequently,
both inequalities of \textup{(S6)} hold.
\end{proposition}

\begin{proof}
At every finite place $w\in\mathrm{supp}$, the local lattice inclusion between $N_\nu$ and
$N_{\nu,0}$ has index determined by the gain $\mathrm{gain}_j(w)$, fixed part $F_j(w)$, and local
energy $E(\nu,j,w)$, with total deviation bounded by $O(X'')$ (\lean{NprimeMuR}). Summing
$\sz(w,\cdot)$ over $w\in\mathrm{supp}$ separates the sum into $\mathcal G(u)+F(u)-E(\nu)$ up to
the local torsion/ramification correction $\mathrm{Lo}(\nu)$. Specifically, \lean{Places_at}
(\lean{TopDensityNGeSixteen.lean:1525}) provides \lean{HypLow}, \lean{HypLow2}, and \lean{HypHigh}
together with $LT\le C n_T m(a/k+b')\log T(\nu)$ and $RX\le C X''$; with $C\ge4$ and
$\mathrm{IOsum}\ge0$ (so $4\,\mathrm{IOsum}\le C\,\mathrm{IOsum}$), feeding these into
Proposition~9.4 of Part~I (\lean{slopesN_at_of}) immediately yields both halves of \textup{(S6)}:
$\mu^+(\nu)\le\mu_0^+(\nu)+\mathcal G(u)+F(u)-E(\nu)+c_1 mB_0$ and
$\mu_0(\nu)+\mathcal G(u)+F(u)-E(\nu)-\mathrm{Lo}(\nu)\le\mu(\nu)$.
\end{proof}

\subsection{Per-place gain inequality (\texttt{LinkG}) and (S7)}

We expand the two core steps behind \lean{LinkG}: extracting half-unit Witt-vector valuations from
Lemma~9.6 of Part~I (\lean{kW_ord}), and summing the orbitwise inequalities into the global
gain-versus-ask bound (\lean{gain_bound_all_of_star}).

\begin{lemma}[Witt-vector order extraction, \texttt{kW\_ord}]\label{lem:kWord}
Let $p$ be a prime, $n'\ge1$, and $O$ be a discrete valuation ring over $W(\mathbb F_{p^{n'}})$
with uniformizer $w_p$ satisfying $w_p^2=p$. Let $g\in\mathrm{GL}_{m+1}(W(\mathbb F_{p^{n'}}))$,
let $s_0,s_1\in\Z$ with $s_1\equiv s_0+d\pmod{n'}$ ($d\ge0$), and put
$M_{b,a}=\sigma^{s_0}(g^{-1})_{b,a}$ and $N'_{b,a}=\sigma^{s_1}(g^{-1})_{b,a}$ in $O$ (where
$\sigma$ is the Witt-vector Frobenius lift). Let $u,u'\in O[X]^{m+1}$ be coordinate vectors under
the scaled lift substitution \lean{liftSubst}, so that $w_p\mid u_b$ and $w_p\mid u'_b$ for all
$b\neq 0$. Then for every $c\in\{0,\dots,m\}$, the polynomial
\[
\Bigl(\sum_{b=0}^m u'_b N'_{b,c}\Bigr)\Bigl(\sum_{b=0}^m u_b M_{b,0}\Bigr)^{p^d}
-\Bigl(\sum_{b=0}^m u'_b N'_{b,0}\Bigr)\Bigl(\sum_{b=0}^m u_b M_{b,c}\Bigr)^{p^d}
\]
is divisible by $w_p$ in $O[X]$ (every coefficient lies in $w_p O$).
\end{lemma}

\begin{proof}
Since $s_1\equiv s_0+d\pmod{n'}$ and $\sigma^{n'}=\mathrm{id}$ on $W(\mathbb F_{p^{n'}})$, the
Frobenius congruence on Witt vectors gives $M_{b,a}^{p^d}\equiv N'_{b,a}\pmod p$ in
$W(\mathbb F_{p^{n'}})$, hence $M_{b,a}^{p^d}=N'_{b,a}+p T_{b,a}$ for some $T_{b,a}\in O$ for all
$a,b$. Because $w_p^2=p$ in $A=O[X]$ and $w_p\mid u_b, u'_b$ for all $b\neq0$, all hypotheses of
Lemma~9.6 of Part~I (\lean{key_minor_dvd}) hold in $A$ with $e=d$. Therefore $w_p$ divides the
twisted $2\times2$ minor in $O[X]$, which means every coefficient lies in $w_p O$.
\end{proof}

\begin{lemma}[Orbitwise arithmetic bounds for cases D and E, \texttt{caseD\_arith}, \texttt{caseE\_arith}]\label{lem:caseDE}
Let $N\ge1$, $n\ge0$, $C_{31}\ge0$, $\ell_x\ge1$, $n\le NG$, and $n_d\le N C_{31}\ell_x$.
\begin{enumerate}[label=(\roman*)]
\item \textup{(Case D: $p<p_0(k)$)} Let $r\ge1$, $b'\ge1$, $0\le c\le u$, $Gu\le b'$,
  $0\le c_0\le 3b'$, $k_Q\le b'$, $0\le D_w\le b'$, $S_1,\mathrm{Pr}_2,t_c\ge0$, and suppose
  $t_c r\le C$ and $2\mathrm{Pr}_2+4t_c+160 C_{31}+7\le C$. Then
  \begin{multline*}
  \mathrm{Pr}_2 r(Nb'+nu)+\bigl(t_c r N(S_1+c_0+k_Q)+2rnc+1\bigr)\\
  +\bigl(40 n_d D_w+2rnc+40 r n_d c_0+1\bigr)+1 \le N(Crb'\ell_x+C\ell_x S_1).
  \end{multline*}
\item \textup{(Case E: $p\ge p_0(k)$)} Let $w,\eta,c,N_m,S_2,E,\mathrm{Pr}_1\ge0$,
  $\mathrm{Pr}_1 G w\le C_w$, $t_{N1}\le C\ell_x$, and
  $\mathrm{Pr}_1 G(\eta+c+N_m)+\kappa_2 B_1+2Gc+40 C_{31}\ell_x E+3\le b_o$. Then
  \begin{multline*}
  \mathrm{Pr}_1 n(w+\eta+c+N_m)+\bigl(t_{N1} N S_2+\kappa_2 N B_1+2nc+1\bigr)\\
  +(40 n_d E+1)+1 \le N(C_w+b_o+C\ell_x S_2).
  \end{multline*}
\end{enumerate}
\end{lemma}

\begin{proof}
For (i), $nu\le NGu\le Nb'$ and $nc\le nu\le Nb'$. Also $n_d D_w\le C_{31}Nb'\ell_x$ and
$n_d c_0\le 3C_{31}Nb'\ell_x$. Since $N,r,b',\ell_x\ge1$, we have $1\le Nb'\le Nb'r\le Nb'r\ell_x$
and $Nb'\ell_x\le Nb'r\ell_x$, while $t_c r N S_1\le C N S_1\le C N\ell_x S_1$. Substituting these
bounds into the left-hand side collects all $Nb'r\ell_x$ coefficients into at most
$(2\mathrm{Pr}_2+4t_c+160 C_{31}+7)Nb'r\ell_x\le C Nb'r\ell_x$, giving (i).

For (ii), substitute $n\le NG$, $n_d\le N C_{31}\ell_x$, $1\le N$, $\mathrm{Pr}_1 Gw\le C_w$, and
$t_{N1}\le C\ell_x$ directly into the left-hand side:
\begin{align*}
\text{LHS} &\le N\bigl(\mathrm{Pr}_1 Gw+\mathrm{Pr}_1 G(\eta+c+N_m)+C\ell_x S_2+\kappa_2 B_1+2Gc+40 C_{31}\ell_x E+3\bigr)\\
&\le N(C_w+b_o+C\ell_x S_2). \qedhere
\end{align*}
\end{proof}

\begin{lemma}[Summation of orbitwise gain bounds, \texttt{gain\_bound\_all\_of\_star}]\label{lem:orbitsum}
Fix an integer $m\ge1$, an inert or signal prime $p\le X''$, a set of orbit representatives
$V\subseteq\mathrm{reps}(p)$ whose top sets $\mathrm{tops}(t_0)$ ($t_0\in V$) are pairwise
disjoint, and nonnegative real quantities $\mathrm{dl}(p,t_0)\ge0$, $\mathrm{ep}(p)\ge0$, and
$g_0(p,j,t_0)$. Suppose that:
\begin{enumerate}[label=(\alph*)]
\item \textup{(Orbitwise bound, \lean{StarAt})} For every $t_0\in V$ and lift $j\in\{1,\dots,m\}$,
  \[
  \sum_{s\in\mathrm{tops}(t_0)}\bigl(\mathrm{ask}_p(j,s)-\mathrm{lo}_G(p)\bigr)\le
  \bigl(g_U(p,j,t_0)-\mathrm{Pr}(p,j,t_0)\bigr)+\mathrm{dl}(p,t_0)-(\lceil\mathrm{tv}(p,t_0)\rceil+D_p+1);
  \]
\item \textup{(Gain and tau definitions)}
  \[
  \mathrm{gain}_p(j)=\sum_{t_0\in V}\max\bigl(g_U-\mathrm{Pr},\max(g_0,0)\bigr)+\sum_{t_0\notin V}\max(g_0,0),
  \]
  and $\tau_p=\mathrm{ep}(p)+m\sum_{t_0\in V}(\lceil\mathrm{tv}(p,t_0)\rceil+D_p+1)$ when
  $V\neq\emptyset$, while $\tau_p=0$ when $V=\emptyset$;
\item \textup{(Floor budget)} $\mathrm{ep}(p)+\sum_{t_0\in\mathrm{reps}(p)}\mathrm{dl}(p,t_0)\le\mathrm{fl}_G(p)$.
\end{enumerate}
Then, with $\mathrm{TTP}(p)=\bigsqcup_{t_0\in V}\mathrm{tops}(t_0)$ the served top set,
\begin{equation}\label{eq:orbitsum}
\sum_{j=1}^m\sum_{s\in\mathrm{TTP}(p)}\bigl(\mathrm{ask}_p(j,s)-\mathrm{lo}_G(p)\bigr)
-m\,\mathrm{fl}_G(p)\le\sum_{j=1}^m\mathrm{gain}_p(j)-\tau_p.
\end{equation}
\end{lemma}

\begin{proof}
When $V=\emptyset$, $\mathrm{TTP}(p)=\emptyset$ (or $p$ is singular / $p\ge X'$, where
$\mathrm{ask}_p\le\mathrm{lo}_G$ or the unserved bound holds) and $\tau_p=0$ by (b), while
$\mathrm{gain}_p(j)\ge0$ and $m\,\mathrm{fl}_G(p)\ge0$, so \eqref{eq:orbitsum} is immediate.

Assume $V\neq\emptyset$. Because the sets $\mathrm{tops}(t_0)$ ($t_0\in V$) are pairwise disjoint,
summing (a) over $t_0\in V$ and then over all $m$ lifts $j\in\{1,\dots,m\}$ yields
\[
\sum_{j=1}^m\sum_{s\in\mathrm{TTP}(p)}\bigl(\mathrm{ask}_p(j,s)-\mathrm{lo}_G(p)\bigr)
\le \sum_{j=1}^m\sum_{t_0\in V}\bigl(g_U-\mathrm{Pr}\bigr)
+m\sum_{t_0\in V}\mathrm{dl}(p,t_0)-m\sum_{t_0\in V}\bigl(\lceil\mathrm{tv}\rceil+D_p+1\bigr).
\]
By (b), $\sum_{t_0\in V}(g_U-\mathrm{Pr})\le\mathrm{gain}_p(j)$ for each $j$ (since the terms
outside $V$ are $\max(g_0,0)\ge0$), and the last term equals $-(\tau_p-\mathrm{ep}(p))$. Since
$m\ge1$ and $\mathrm{ep}(p),\mathrm{dl}(p,t_0)\ge0$, (c) gives
\[
\mathrm{ep}(p)+m\sum_{t_0\in V}\mathrm{dl}(p,t_0)
\le m\Bigl(\mathrm{ep}(p)+\sum_{t_0\in\mathrm{reps}(p)}\mathrm{dl}(p,t_0)\Bigr)\le m\,\mathrm{fl}_G(p).
\]
Rearranging gives \eqref{eq:orbitsum}.
\end{proof}

\begin{theorem}[Per-place gain inequality and verification of (S7)]\label{thm:S7}
Assume \ref{C8}. For every inert or signal place $w\mid p\le X''$, the gain clause
\eqref{eq:orbitsum} of \lean{LinkG_at} holds. Weighting \eqref{eq:orbitsum} by $\sz(w,\cdot)$ and
summing over $w\mid p\le X''$ via the product formula $\sum_{w\mid p}\mathrm{lw}(w)e_w=\log p$ gives
\[
A(u)-\mathcal G(u)\le L_P=c_L n_T m b\,x^{e_L\eps},\qquad I(u)\le L_P,
\]
establishing \textup{(S7)}.
\end{theorem}

\begin{proof}
Lemma~\ref{lem:kWord} (\lean{kW_ord}) together with \lean{kS_ord} establishes
\lean{CoarseMultiplicityOrderBound} (\lean{key3Ord}), supplying a half-unit of $p$-adic valuation
at every odd exponent $p^e$ in the prime-power sum
\[
\mathrm{vonMangoldtClassSum}_{j,t}(p)=\sum_{i\text{ odd}}\max\bigl(0,\,a_t/p^i-w_{\mathrm{part}}(j,t,p,i)\bigr).
\]
Splitting at $p_0(k)$ (cases D and E via Lemma~\ref{lem:caseDE} in \lean{St.starAll_of})
establishes the orbitwise bound (a) of Lemma~\ref{lem:orbitsum}, and Lemma~\ref{lem:orbitsum}
produces \eqref{eq:orbitsum}.

Recall from Part~I that
$\ASK=\sum_{j,t}\bigl(\sum_{p\text{ inert}}2\log p\cdot\mathrm{vonMangoldtClassSum}_{j,t}(p)
+\log d\cdot\mathrm{ask}_D\bigr)$.
Multiplying \eqref{eq:orbitsum} by $2\mathrm{lw}(w)e_w$ and summing over all places $w\mid p$
uses $\sum_{w\mid p}\mathrm{lw}(w)e_w=\log p$. The error terms $\mathrm{lo}_G(p)$ and
$\mathrm{fl}_G(p)$ (recorded in \lean{GlobalTheoremInputs.lean}) sum over $p\le X''$ to at most
$c_L n_T m b\,x^{e_L\eps}=L_P$: in particular, the additive term $b_o=b^{\delta(x)}$ inside
$\mathrm{lo}_G(p)$ satisfies $b^{\delta(x)}\le b^{1/2}\le b$ once $\delta(x)\le\frac12$ (since
$\delta(x)\to0$ as $x\to\infty$), so its sum over $j,t,p$ is absorbed into the factor $b$ of
$L_P=c_L n_T m b\,x^{e_L\eps}$. Similarly, the local inclusion length satisfies $I(u)\le L_P$
(\lean{Loss.loss_bound_at}), proving \textup{(S7)}.
\end{proof}

% ======================================================================
\section{Near-orthogonality (\texttt{LinkS}), the \texorpdfstring{$\Sha$}{Sha}-window, and (S8)--(S10)}\label{sec:S8VWin}
% ======================================================================

\subsection{Near-orthogonality and the frame inequality (S8)--(S10)}

We first record the elementary non-degeneracy consequence of the diagonal Hilbert--Schmidt bound
(\lean{unfolding_constant_pos_of_diagonal}, \lean{dall_pos} in \lean{TopDensityNGeSixteen.lean}).

\begin{lemma}[Non-degeneracy from the diagonal Gram bound]\label{lem:diagnondeg}
Let $C_\nu$ and $\mathcal H$ be finite-dimensional complex inner product spaces of dimensions
$d(\nu)=\dim C_\nu>0$ and $D_{\mathrm{all}}=\dim\mathcal H$, let $P_\nu:C_\nu\to\mathcal H$ be
linear, let $c_0(\nu)\ge0$, and put $K(\nu,\nu)=c_0(\nu)^{1/2}P_\nu^\ast P_\nu$. If
$\|K(\nu,\nu)-I_{C_\nu}\|_{HS}^2\le e_0 d(\nu)$ with $e_0<1$, then $c_0(\nu)>0$ and
$D_{\mathrm{all}}>0$.
\end{lemma}

\begin{proof}
In an orthonormal basis of $C_\nu$, $\|0-I_{C_\nu}\|_{HS}^2=\sum_{\ell=1}^{d(\nu)}1=d(\nu)$. If
either $c_0(\nu)=0$ or $D_{\mathrm{all}}=0$ (so $\mathcal H=\{0\}$ and $P_\nu=0$), then
$K(\nu,\nu)=0$, giving $d(\nu)=\|K(\nu,\nu)-I_{C_\nu}\|_{HS}^2\le e_0 d(\nu)<d(\nu)$, a
contradiction.
\end{proof}

\begin{lemma}[Genus-mass comparability proviso, \texttt{mass\_corollary}, \texttt{proviso\_e}]\label{lem:proviso}
Let $f,k\in\N$ with $1\le f\le k$, let $T\ge2$, $\Gamma>1$, $D^\ast>0$, and $d^\ast>0$, and set
\[
X=\Bigl(\frac{D^\ast}{d^\ast}\Bigr)^{1/(fk)},\qquad
Y=T^\ast=X\,\Gamma^{-2/(fk)}.
\]
Suppose that the word prime $T=T(\nu)$ lies in the window
\begin{equation}\label{eq:Twin}
Y\le T\le Y\Bigl(1+\frac{1}{fk}\Bigr),
\end{equation}
and that the genus-mass rank identity $d(\nu)=d^\ast(T^k-1)^f$ holds. Then:
\begin{enumerate}[label=(\roman*)]
\item $\frac12\,D^\ast\Gamma^{-2}\le d(\nu)\le e\,D^\ast\Gamma^{-2}$;
\item $\log X-\frac{2\log\Gamma}{fk}\le\log T\le\log 2+\log X$;
\item in particular, $d(\nu)>0$ and the comparability proviso \lean{ChosenProviso} holds with
  explicit constant $C_0=e$:
  \begin{equation}\label{eq:proviso_e}
  e^{-1}\le\frac{d(\nu)\,\Gamma^2}{D^\ast}\le e.
  \end{equation}
\end{enumerate}
\end{lemma}

\begin{proof}
By definition of $X$ and $Y$, we have $X^{fk}=D^\ast/d^\ast$ and hence
$d^\ast Y^{fk}=d^\ast X^{fk}\Gamma^{-2}=D^\ast\Gamma^{-2}$. Put $N=fk\ge1$.

For the upper bound in (i), $1+1/N\le e^{1/N}$ gives $(1+1/N)^N\le e$
(\lean{one_add_inv_pow_le_exp_one}). Using $T^k-1\le T^k$ and the upper half of \eqref{eq:Twin},
\[
(T^k-1)^f\le T^{fk}\le Y^{fk}\Bigl(1+\frac{1}{fk}\Bigr)^{fk}\le e\,Y^{fk}.
\]
Multiplying by $d^\ast>0$ and substituting $d(\nu)=d^\ast(T^k-1)^f$ and
$d^\ast Y^{fk}=D^\ast\Gamma^{-2}$ yields $d(\nu)\le e\,D^\ast\Gamma^{-2}$.

For the lower bound in (i), $1\le f\le k$ and $k-1<2^{k-1}$ imply $2f\le 2k\le 2\cdot 2^{k-1}=2^k$.
Since $T\ge2$, we have $2f\le 2^k\le T^k$, so $f/T^k\le 1/2$ and $1/T^k\le 1$. By Bernoulli's
inequality $(1-z)^f\ge 1-fz$ for $0\le z\le 1$,
\[
\Bigl(1-\frac{1}{T^k}\Bigr)^f\ge 1-\frac{f}{T^k}\ge\frac12.
\]
Combining this with $Y^{fk}\le T^{fk}$ from the lower half of \eqref{eq:Twin},
\[
(T^k-1)^f=T^{fk}\Bigl(1-\frac{1}{T^k}\Bigr)^f\ge\frac12\,T^{fk}\ge\frac12\,Y^{fk}.
\]
Multiplying by $d^\ast$ gives $d(\nu)\ge\frac12\,D^\ast\Gamma^{-2}$, completing (i).

For (ii), since $\Gamma>1$ and $N\ge1$ we have $\Gamma^{-2/N}\le 1$, so $Y\le X$ and hence
$T\le Y(1+1/N)\le 2Y\le 2X$. Taking logarithms in $Y\le T\le 2X$ yields (ii).

Finally, for (iii), $e\ge 1+1=2$ gives $e^{-1}\le 1/2$. Multiplying (i) by $\Gamma^2/D^\ast>0$
therefore yields \eqref{eq:proviso_e}.
\end{proof}

\noindent In the live Lean closure, the genus-mass rank identity $d(\nu)=d^\ast(T^k-1)^f$
(\lean{LinkB.RankM}) is supplied by \ref{C8} from three ingredients: the Hermitian class-number
formula \lean{ClassNumberFormula}$_2$ (\lean{classNumber_closed}), adelic quotient volume
independence of signature \lean{VolSigIndep} (\lean{volSigIndep_and_classItemsAt_of_massRatioUnderVolume}),
and the splitting/level-tie hypotheses \lean{RankTiesF} (\lean{rank_conditions_iff_lowest_space}:
the host primes $g_0,g_1$ and word primes $T(\nu)$ split in the totally real subfield $F^+$,
together with \lean{Mu0Tie}). Lemma~\ref{lem:proviso} then closes \lean{ChosenProviso} completely
without any further automorphic input.

\begin{lemma}[Irreducibility decay, \texttt{delp\_of\_allRows}, \texttt{linkIRR\_of\_allRows}]\label{lem:linkIRR}
Let $0<\eps_z\le 1$, $C_s\ge 6$, and $c>0$. At scale $x\ge\max(10, 100 C_s^2+1, 1/c)$, let
$f,m,N,g_0,g_1,q,d\in\N$ satisfy the size regime \lean{SizesAt}:
\begin{gather*}
g_k\ge\exp(8 x^{\eps_z})\ (k\in\{0,1\}),\qquad \min(q,d)\ge\exp(x^{\eps_z/6}),\\
x^2\le f\le x^{C_s},\qquad m\le x^{C_s},\qquad 100\le N\le x^{C_s}.
\end{gather*}
For $k\in\{0,1\}$, define the local bad-type weight
\begin{multline}\label{eq:Wg}
W_g(f,m,N,g_k,q,d)=\frac{3f}{g_k}+\frac{23f\log g_k}{\sqrt{g_k}}
+2fm\Bigl(\frac{1}{4\min(q,d)}+\frac{1+\log N}{g_k-1}\Bigr)\\
+\frac{N}{2}\Bigl(\frac{7}{8}\Bigr)^f+\exp\Bigl(-\frac{f}{5\log N}\Bigr),
\end{multline}
and put $\varepsilon_C=\exp(-c N f\log N)$. Suppose that the total row count $D_{\mathrm{all}}$,
the good (residually irreducible) row count $D_{\mathrm{rows}}\le D_{\mathrm{all}}$ (with
$\mathrm{del}_p\cdot D_{\mathrm{all}}=D_{\mathrm{all}}-D_{\mathrm{rows}}$), the bad-row counts
$B_0,B_1$ at $g_0,g_1$, the main term $d_S>0$, and the local bad-cylinder shares $p_0,p_1\ge0$
satisfy:
\begin{enumerate}[label=(\alph*)]
\item \textup{(Union bound)} $D_{\mathrm{all}}-D_{\mathrm{rows}}\le B_0+B_1$;
\item \textup{(Cylinder count, \lean{CountAllRows})} $|D_{\mathrm{all}}-d_S|\le\varepsilon_C d_S$
  and $|B_k-p_k d_S|\le\varepsilon_C p_k d_S$ for $k\in\{0,1\}$;
\item \textup{(Local bad-share bound, \lean{WeightAllRows})}
  $p_k\le W_g(f,m,N,g_k,q,d)$ for $k\in\{0,1\}$.
\end{enumerate}
Then:
\begin{enumerate}[label=(\roman*)]
\item $\varepsilon_C\le\frac12$ and $W_g(f,m,N,g_k,q,d)\le x^{3C_s+3}\exp(-x^{\eps_z/6})$ for
  $k\in\{0,1\}$;
\item $0\le\mathrm{del}_p\le \frac{1+\varepsilon_C}{1-\varepsilon_C}(W_0+W_1)\le
  6\,x^{3C_s+3}\exp(-x^{\eps_z/6})$;
\item for all sufficiently large $x\ge x_0(\eps_z,C_s,c)$, we have
  $0\le\mathrm{del}_p\le\exp(-x^{\eps_z/7})$, establishing \lean{LinkIRR}.
\end{enumerate}
\end{lemma}

\begin{proof}
Put $P=x^{C_s}\ge 100$, $E=x^{\eps_z}\ge 1$, $u=x^{\eps_z/6}\in[1,x]$, and $e_u=\exp(-u)$. Since
$x\ge 10$, we have $x^3\ge 1000\ge 103$.

For $\varepsilon_C\le\frac12$ (\lean{epsilon_bound_large_x}), $x\ge\max(1,1/c)$ gives
$c f\ge c x^2\ge c x\ge 1$, and $N\ge 100>e$ gives $N\log N\ge 1$. Hence
$c N f\log N\ge 1$, so $\varepsilon_C\le e^{-1}\le\frac12$.

For the bound on $W_k=W_g(f,m,N,g_k,q,d)$ (\lean{weight_numbers_bound}), fix $g=g_k$. Since
$g\ge\exp(8E)\ge 1+8E\ge 9$ and $8E\ge u$, we have $e_u g\ge\exp(8E-u)\ge 1$, so $1/g\le e_u$. We
bound the five terms of \eqref{eq:Wg} in turn:
\begin{enumerate}[label=\arabic*.]
\item $3f/g\le 3P e_u$ since $f\le P$ and $1/g\le e_u$.
\item Set $s=g^{1/4}\ge\exp(2E)\ge e^u$, so $1/s\le e_u$. Using $\log s\le s-1\le s$, we get
  $\log g=4\log s\le 4s$, whence $23f\log g/\sqrt{g}=23f\log g/s^2\le 92f/s\le 92P e_u$.
\item Since $\min(q,d)\ge e^u$, we have $\frac{1}{4\min(q,d)}\le e_u$. Moreover, $e_u g\ge 1$ and
  $g\ge 9$ give $e_u(g-1)\ge\frac{8}{9}>\frac12$, while $1+\log N\le N\le P$, so
  $(1+\log N)/(g-1)\le 2P e_u$. Since $f,m\le P$, the third term is at most
  $2P^2(e_u+2P e_u)=2P^2 e_u+4P^3 e_u$.
\item Since $7/8\le e^{-1/8}$ and $f\ge x^2\ge 8x\ge 8u$ (as $x\ge 10$), we have
  $(7/8)^f\le e^{-f/8}\le e_u$, so $(N/2)(7/8)^f\le P e_u$.
\item Since $N\le x^{C_s}$ and $\log x=2\log(\sqrt x)\le 2\sqrt x$, we have
  $\log N\le C_s\log x\le 2 C_s\sqrt x$. Because $x\ge 100 C_s^2+1$, we have $10 C_s\le\sqrt x$,
  hence $5\log N\le 10 C_s\sqrt x\le x$ and $u\cdot 5\log N\le x^2\le f$. Thus
  $\exp(-f/(5\log N))\le e_u$.
\end{enumerate}
Summing these five bounds and using $1\le P\le P^2\le P^3$ and $103\le x^3$ (\lean{wsum}) gives
\[
W_k\le (3P+92P+2P^2+4P^3+P+1)e_u\le 103 P^3 e_u\le P^3 x^3 e_u=x^{3C_s+3}\exp(-x^{\eps_z/6}),
\]
proving (i).

For (ii) (\lean{delp_of_allRows}, \lean{share_from_count}), hypothesis~(b) gives
$D_{\mathrm{all}}\ge(1-\varepsilon_C)d_S>0$ and $B_k\le(1+\varepsilon_C)p_k d_S\le(1+\varepsilon_C)W_k d_S$
by (c). Since $D_{\mathrm{rows}}\le D_{\mathrm{all}}$, we have $\mathrm{del}_p\ge 0$, and by (a),
\[
\mathrm{del}_p=\frac{D_{\mathrm{all}}-D_{\mathrm{rows}}}{D_{\mathrm{all}}}
\le\frac{B_0+B_1}{D_{\mathrm{all}}}
\le\frac{1+\varepsilon_C}{1-\varepsilon_C}(W_0+W_1).
\]
Since $\varepsilon_C\le\frac12$, $\frac{1+\varepsilon_C}{1-\varepsilon_C}\le 3$
(\lean{numbers_dominant_large}), so (i) yields $\mathrm{del}_p\le 6\,x^{3C_s+3}\exp(-x^{\eps_z/6})$.

Finally, for (iii), since $\eps_z/6>\eps_z/7$, we have $6\,x^{3C_s+3}\exp(x^{\eps_z/7}-x^{\eps_z/6})\to 0$
as $x\to\infty$ (\lean{dominant_numbers_sharp_holds}), giving $\mathrm{del}_p\le\exp(-x^{\eps_z/7})$
for $x\ge x_0(\eps_z,C_s,c)$.
\end{proof}

\noindent In the live Lean closure, \lean{sizesRec_at} (\lean{LinkIRROfJoinEx.lean:83}) verifies
\lean{SizesAt} with $C_s=6$ directly from the scale windows of Section~\ref{sec:host}, while
\lean{locShare_le_Wg} (\lean{ConcNWeight.lean:284}) establishes hypothesis~(c) from Green's
cuspidal parametrisation (\lean{admLq}, \lean{GlobalTheory.GLq.Green}), comparison transfer
($(1-1/g_k)P(c)\le w_J(v,c)$), and the finite-field polynomial irreducibility estimate
(\lean{property_W_of_estimates}). Combined with the cylinder count \lean{CountAllRows} from
\ref{C8}, Lemma~\ref{lem:linkIRR} closes \lean{LinkIRR} (\lean{closed_linkIRR_of_join_count},
\lean{LinkIrr.lean:60}).

\begin{lemma}[Gram decorrelation assembly, \texttt{linkS\_concl}, \texttt{eF\_nonneg\_and\_le\_exp\_neg}]\label{lem:linkS}
Let $C_0\ge 1$, $C_\Gamma>0$, $c_7>0$, $0<\eps<1$, and put $y=x^{\eps/8}$, $L=x^\eps=y^8$,
$b=\exp(L)$, $\Lambda=C_\Gamma m^5 f L$, and $\Gamma=\exp(\Lambda)$. Suppose that $m\ge 2$,
$f\ge 1$, $r\ge k\ge 5$, and $k\ge 8\Lambda$, and that the comparability proviso
$d(\nu)\Gamma^2/D^\ast\le C_0$ holds for every $\nu\in\mathcal W(u)$. Suppose further that the
Gram operator $K(\nu,\nu')$ decomposes via its degenerate (parabolic) part $K_{\mathrm{pp}}(\nu,\nu')$
so that:
\begin{enumerate}[label=(\alph*)]
\item \textup{(\texttt{ClKppOff})} $K_{\mathrm{pp}}(\nu,\nu')=0$ for all $\nu\neq\nu'$;
\item \textup{(\texttt{ClKppDiag})} $\|K_{\mathrm{pp}}(\nu,\nu)\|_{HS}^2\le d(\nu)\exp(-kf\log k)$
  for all $\nu$;
\item \textup{(\texttt{ClOff})} $\|K(\nu,\nu')-K_{\mathrm{pp}}(\nu,\nu')\|_{HS}^2\le
  (1+b^{-6})\,d(\nu)d(\nu')/D^\ast$ for all $\nu\neq\nu'$;
\item \textup{(\texttt{ClDiag})} $\|K(\nu,\nu)-K_{\mathrm{pp}}(\nu,\nu)-I_{C_\nu}\|_{HS}^2\le
  \Gamma\,d(\nu)^2/D^\ast$ for all $\nu$;
\item \textup{(\texttt{adm7})} $|D_{\mathrm{all}}-D^\ast|\le\exp(-c_7 r f\log r)\,D^\ast$.
\end{enumerate}
Set $\eta_0=b^{-6}$, $\eta_1=\exp(-c_7 r f\log r)$, and
$e_0=2C_0/\Gamma+2\exp(-kf\log k)$. Then:
\begin{enumerate}[label=(\roman*)]
\item for all $\nu\neq\nu'$, $\|K(\nu,\nu')\|_{HS}^2\le(1+\eta_0)\,d(\nu)d(\nu')/D^\ast$;
\item for all $\nu$, $\|K(\nu,\nu)-I_{C_\nu}\|_{HS}^2\le e_0\,d(\nu)$, and
  $D_{\mathrm{all}}\le(1+\eta_1)D^\ast$;
\item for all $y\ge Y_0=\max\bigl(1,\,\frac{2+\log(2C_0+2)}{C_\Gamma},\,\frac{1}{8c_7 C_\Gamma}\bigr)$
  (i.e. $x\ge Y_0^{8/\eps}$), we have $e_0\ge 0$ and
  \[
  \eta_0,\ \eta_1,\ \sqrt{e_0}\le\exp(-y)=\exp(-x^{\eps/8}).
  \]
\end{enumerate}
\end{lemma}

\begin{proof}
Part (i) is immediate from (a) and (c), and $D_{\mathrm{all}}\le(1+\eta_1)D^\ast$ is immediate
from (e) (\lean{linkS_concl}).

For the diagonal bound in (ii) (\lean{diag_unproj}), write
$K(\nu,\nu)-I_{C_\nu}=(K(\nu,\nu)-K_{\mathrm{pp}}(\nu,\nu)-I_{C_\nu})+K_{\mathrm{pp}}(\nu,\nu)$.
In an orthonormal basis $\{v_\ell\}$ of $C_\nu$,
$\|(A+B)v_\ell\|^2\le 2\|Av_\ell\|^2+2\|Bv_\ell\|^2$, so summing over $\ell$ gives
$\|A+B\|_{HS}^2\le 2\|A\|_{HS}^2+2\|B\|_{HS}^2$ (\lean{hsNormSq_add_le}). Applying this with (d)
and (b) yields
\[
\|K(\nu,\nu)-I_{C_\nu}\|_{HS}^2
\le 2\Gamma\,\frac{d(\nu)^2}{D^\ast}+2\exp(-kf\log k)\,d(\nu).
\]
By the comparability proviso $d(\nu)/D^\ast\le C_0\Gamma^{-2}$, we have
$\Gamma\,d(\nu)^2/D^\ast\le (C_0/\Gamma)\,d(\nu)$, giving $\|K(\nu,\nu)-I_{C_\nu}\|_{HS}^2\le e_0\,d(\nu)$.

For (iii) (\lean{size_thresholds}, \lean{eF_nonneg_and_le_exp_neg}), since $m\ge 2$ and $f\ge 1$,
we have $\Lambda=C_\Gamma m^5 f y^8\ge C_\Gamma y^8$. Since $r\ge k\ge 5>e$, we have
$\log k\ge 1$ and $\log r\ge 1$, so
\[
r f\log r\ge k f\log k\ge k\ge 8\Lambda\ge 8 C_\Gamma y^8.
\]
Now assume $y\ge Y_0$, so $y\ge 1$ and $y^2\le y^8$. Then:
\begin{itemize}
\item $\eta_0=\exp(-6y^8)\le\exp(-y)$ since $6y^8\ge y$;
\item $c_7 r f\log r\ge 8 c_7 C_\Gamma y^8\ge 8 c_7 C_\Gamma y^2\ge y$ (since
  $y\ge(8c_7 C_\Gamma)^{-1}$), so $\eta_1=\exp(-c_7 r f\log r)\le\exp(-y)$;
\item both $\Gamma^{-1}=e^{-\Lambda}$ and $e^{-kf\log k}$ are at most $\exp(-C_\Gamma y^8)$, and
  $C_\Gamma y^8\ge C_\Gamma y^2\ge(2+\log(2C_0+2))y\ge 2y+\log(2C_0+2)$, whence
  \[
  0\le e_0\le(2C_0+2)\exp(-C_\Gamma y^8)=\exp\bigl(\log(2C_0+2)-C_\Gamma y^8\bigr)\le\exp(-2y).
  \]
  Taking square roots gives $\sqrt{e_0}\le\exp(-y)=\exp(-x^{\eps/8})$. \qedhere
\end{itemize}
\end{proof}

\noindent In the live Lean closure, the four kernel clauses (a)--(d) and the point-counting bound
(e) are supplied by \ref{C8} via \lean{KppOffC_holds}, \lean{kernel_diagonal_positive}
(\lean{RowsPos.lean:221}), \lean{diagonal_positive_final} (\lean{RowsPos.lean:261}), and
\lean{CountAllPointsClass} (\lean{countOwn_of_rowsPEv}, \lean{LinkSFinal.lean:57}); Lemma~\ref{lem:linkS}
(\lean{linkS_of_standing_hypotheses}, \lean{StandingHypothesesHoldFamily.lean:598}) then closes
\lean{LinkS} inside \lean{theoremS_link_concrete} (\lean{LinkSPos.lean:151}).

\begin{proposition}[Verification of \texttt{LinkS} and (S8)--(S10)]\label{prop:S8S10}
Assume \ref{C8} and put $\Gamma=\exp(C_\Gamma m^5 f L)$. Then for $x\ge x_0(\eps)$:
\begin{enumerate}[label=(\roman*)]
\item \textup{(Comparability proviso)} Every word $\nu\in\mathcal W(u)$ satisfies
  $e^{-1}\le d(\nu)\Gamma^2/D^\ast\le e$ (\lean{ChosenProviso} with $C_0=e$, by
  Lemma~\ref{lem:proviso}); in particular $d(\nu)>0$.
\item \textup{(Four Gram clauses and \lean{LinkS})} By Lemma~\ref{lem:linkS} (from the four kernel
  clauses \lean{ClKppOff}, \lean{ClKppDiag}, \lean{ClOff}, \lean{ClDiag}, residual evaluation
  \lean{ResidEv}, and point counting \lean{CountAllPointsClass}), there exist
  $\eta_0,\eta_1,e_0\ge0$ with $\eta_0,\eta_1,\sqrt{e_0}\le\exp(-x^{\eps/8})$ such that
  \begin{gather*}
  \|K(\nu,\nu')\|_{HS}^2\le(1+\eta_0)\frac{d(\nu)d(\nu')}{D^\ast}\quad(\nu\neq\nu'),\\
  \|K(\nu,\nu)-I\|_{HS}^2\le e_0\,d(\nu),\qquad D_{\mathrm{all}}\le(1+\eta_1)D^\ast.
  \end{gather*}
\item \textup{(Irreducibility decay, \lean{LinkIRR})} By Lemma~\ref{lem:linkIRR}, the
  irreducibility error satisfies $0\le\delta_u=\mathrm{del}_p\le\exp(-x^{\eps/7})$.
\item Feeding (i)--(iii) and Lemma~\ref{lem:diagnondeg} into Theorem~9.2 of Part~I (Frame theorem)
  yields \textup{(S8)}, \textup{(S9)}, and \textup{(S10)} with
  $(\bar\eta+4\sqrt{\bar e})x^{c_{HH}}\le 6$.
\end{enumerate}
\end{proposition}

\begin{proof}
Part (i) follows directly from Lemma~\ref{lem:proviso} (\lean{proviso_e}, yielding
\lean{chosenProviso_4f} from \lean{LinkB.RankM} and the word prime window \lean{D.T_win}), part
(ii) follows from Lemma~\ref{lem:linkS} (\lean{linkS_of_standing_hypotheses},
\lean{theoremS_link_concrete}, \lean{LinkSPos.lean:151}), and part (iii) follows from
Lemma~\ref{lem:linkIRR} (\lean{linkIRR_of_allRows}, \lean{closed_linkIRR_of_join_count}, yielding
\lean{LinkIRR} of \lean{GlobalTheoremInputs.lean:64}). For $x\ge1$ we have $\sqrt{e_0}\le\exp(-x^{\eps/8})<1$, so
$e_0<1$; since $d(\nu)>0$ by (i), Lemma~\ref{lem:diagnondeg} gives $c_0(\nu)>0$ and
$D_{\mathrm{all}}>0$ (\lean{matKF_of}, \lean{rrRankF_of}). For $x\ge x_0(\eps)$, part~(iii) gives
$\delta_u\le\exp(-x^{\eps/7})<1/2$, while part~(ii) gives $0\le e_u\le\frac14$ and
$(\eta_u+4\sqrt{e_u})x^{c_{HH}}\le 6$ (since $\exp(-x^{\eps/8})$ decays faster than any power of
$x$). The height cap $0<H_\ast(u)\le x^{c_{HH}}b$ follows from
$s(u)-\mu_\ast^-(u)\le s(u)+\mu_{\max}(R)\le x^{c_{HH}}b$ (\lean{LinkDIV_at}, second clause).
Applying Theorem~9.2 of Part~I (as noted in Remark~9.3 of Part~I) immediately gives \textup{(S8)},
\textup{(S9)}, and \textup{(S10)}.
\end{proof}

\subsection{The Tate--Shafarevich window \texttt{VWin}}

The only remaining ingredient of Theorem~C2 is the existence of arbitrarily large scale windows
$[Y,Y^2]$ (which contain $[Y,Y^{3/2}]$ for Theorem~A of Part~I) on which the cyclotomic
Tate--Shafarevich loss is at most $\sqrt{b}=\exp(x^\eps/2)$. We isolate the explicit minus-part
calculation (\lean{ThmBcOfPsiBound.lean:220--340}) as a standalone lemma.

\begin{lemma}[Explicit minus-part bound, \texttt{ThmBcFree}]\label{lem:minuspart}
Let $E$ be an abelian number field of degree $d_E=[E:\Q]$ and conductor $F\ge1$, and let $t\ge2$.
Suppose that:
\begin{enumerate}[label=(\alph*)]
\item \textup{(Primes $\ell\nmid d_E$, \lean{PrimeBound})} for every $1\le e'<t$ and finite set
  $S'$ of odd primes not dividing $d_E$,
  $\sum_{\ell\in S'}\log\ell\cdot\len\Sha_\ell^-(e')\le d_E(2t+5)\log(tF)$;
\item \textup{(Primes $\ell\mid d_E$, \lean{StickBound})} there is $C\ge0$ such that for every odd
  prime $\ell\mid d_E$ and $1\le e'<t$,
  $\len\Sha_\ell^-(e')\log\ell\le C\ell d_E^2(e'+1)\log(e'\ell F)^2$.
\end{enumerate}
Then for every finite set $S$ of odd primes,
\begin{equation}\label{eq:minusbound}
\sum_{\ell\in S}\log\ell\sum_{1\le e'<t}\len\Sha_\ell^-(e')
\le 3\,d_E t^2\log(tF)+C\sum_{\substack{\ell\mid d_E\\\ell\text{ odd}}}\ell\,d_E^2 t^2\log(t\ell F)^2.
\end{equation}
\end{lemma}

\begin{proof}
Partition $S$ into $S_1=\{\ell\in S:\ell\nmid d_E\}$ and $S_2=\{\ell\in S:\ell\mid d_E\}$.
For $S_1$, swapping the order of summation and applying (a) for each of the $t-1$ values
$e'\in\{1,\dots,t-1\}$ gives
\[
\sum_{e'=1}^{t-1}\sum_{\ell\in S_1}\log\ell\cdot\len\Sha_\ell^-(e')
\le (t-1)\,d_E(2t+5)\log(tF)\le 3\,d_E t^2\log(tF),
\]
using $(t-1)(2t+5)=2t^2+3t-5\le 3t^2$ for $t\ge2$.

For each $\ell\in S_2$ and $1\le e'<t$, we have $e'+1\le t$ and $1\le e'\ell F\le t\ell F$, so
$\log(e'\ell F)^2\le\log(t\ell F)^2$. Applying (b) and summing over $1\le e'<t$ yields
\[
\log\ell\sum_{e'=1}^{t-1}\len\Sha_\ell^-(e')
\le (t-1)\cdot C\ell d_E^2 t\log(t\ell F)^2\le C\ell d_E^2 t^2\log(t\ell F)^2.
\]
Summing over $\ell\in S_2$ (which is a subset of the odd prime factors of $d_E$) and adding the
$S_1$ bound gives \eqref{eq:minusbound}.
\end{proof}

\begin{lemma}[Combining bad and good primes, \texttt{point\_arith}]\label{lem:pointHV}
Let $B_\psi$ be a finite set of primes of cardinality $K=|B_\psi|$, let $M=Nt\ge2$, $B\ge0$,
$C_0\ge0$, and suppose:
\begin{enumerate}[label=(\alph*)]
\item \textup{(\texttt{SumOut})} for every finite set $S'$ of odd primes disjoint from $B_\psi$,
  \[
  \sum_{\ell\in S'}\log\ell\sum_{1\le e'<t}\len\Sha_\ell^+(e')\le M^{C_0};
  \]
\item \textup{(\texttt{SumIn})} for each prime $\ell\in B_\psi$, there is $C_\ell\in\N$ such that
  $\log\ell\sum_{1\le e'<t}\len\Sha_\ell^+(e')\le M^{C_\ell}+C_\ell B$.
\end{enumerate}
Put $D=\sum_{\ell\in B_\psi}C_\ell$ and $C_\ast=C_0+D+K+KD$. Then for every finite set $S$ of odd
primes,
\begin{equation}\label{eq:pointHV}
\sum_{\ell\in S}\log\ell\sum_{1\le e'<t}\len\Sha_\ell^+(e')\le M^{C_\ast}+C_\ast B.
\end{equation}
\end{lemma}

\begin{proof}
Split $S$ into $S_{\mathrm{in}}=S\cap B_\psi$ (of cardinality $k\le K$) and
$S_{\mathrm{out}}=S\setminus B_\psi$. By (a), the sum over $S_{\mathrm{out}}$ is at most $M^{C_0}$.
For each $\ell\in S_{\mathrm{in}}$, $C_\ell\le D$ and $M\ge2>1$, so (b) gives at most $M^D+DB$,
and summing over $S_{\mathrm{in}}$ gives at most $k(M^D+DB)$. Since $M\ge2$, Bernoulli's inequality
yields $M^K=(1+(M-1))^K\ge 1+K(M-1)\ge K+1$. Since $C_0\ge0$ gives $M^{C_0}\ge1$, and
$M^D,M^{KD}\ge1$ with $kD\le KD\le C_\ast$, we obtain
\[
k(M^D+DB)+M^{C_0}\le (K+1)M^{C_0+D}+kDB\le M^{C_0+D+K+KD}+C_\ast B=M^{C_\ast}+C_\ast B. \qedhere
\]
\end{proof}

\begin{theorem}[The $\Sha$-window \texttt{VWin}]\label{thm:VWin}
Assume \ref{C4}, \ref{C5}, and \ref{C7}. Let $\psi$ be a primitive Dirichlet character of order
$o$, $C_t=401$, $0<\eps<1/52$, and $e(x)=\lceil(2+8\eps)\log_2 x\rceil$. Then there exist $J\ge2$
and $X_0>0$ such that for every $X\ge X_0$ and every window pair $(q,d)$ at $(\eps,X^{1/J})$,
there exists $Y$ with $X^{1/J}\le Y$ and $Y^2\le X$ such that for every $x\in[Y,Y^2]$, the abelian
field $E$ of Section~\ref{sec:glued} (\lean{IsChiField}) satisfies
\begin{equation}\label{eq:VVle}
\sum_{\ell\text{ prime}}\log\ell\sum_{1\le e'<\lfloor x^{401}\rfloor}\len\Sha_\ell(e')
\le \sqrt{\exp(x^\eps)}=\sqrt{b}.
\end{equation}
\end{theorem}

\begin{proof}
We split $\Sha_\ell(e')$ into its minus part $\Sha_\ell^-(e')$ and plus part $\Sha_\ell^+(e')$
under complex conjugation (\lean{vwin_of_counts}, \lean{JointWinOfCounts.lean:73}); the prime
$\ell=2$ (both minus and plus parts) is absorbed into the $\ell=2$ Iwasawa tower bound
(\lean{Two.twoIn}) in Step~2 below, so Step~1 treats all odd primes $\ell$.

\emph{Step 1: The odd-prime minus part (\lean{ThmBcFree}).}
By Poitou--Tate duality and unramified orthogonality \ref{C5} (\lean{primeBound_of}) and
Stickelberger annihilators (\lean{stickBound}), hypotheses (a)--(b) of Lemma~\ref{lem:minuspart}
hold for $E$ with $d_E=[E:\Q]\le 4o\cdot 2^{e(x)}\ll_\psi x^{2+8\eps}$,
$\log F\ll x^{\eps/5}(\log x)^2$, and $t=\lfloor x^{401}\rfloor$. By \eqref{eq:minusbound}, the
entire odd-prime minus-part sum is bounded by a polynomial in $x$ and $\log F$, hence is
$o(\exp(x^\eps/2))$ for $x\ge X_0$.

\emph{Step 2: The plus part and $\ell=2$ (\lean{ThmBab} and \lean{JointWinOfCounts}).}
By \ref{C4} (\lean{UpperTheory.V.FW.ferreroWashington}) and Poitou--Tate duality \ref{C5}
(\lean{PropA.aBound_of_isTower}), the fixed-prime and $\ell=2$ Iwasawa tower bounds hold
(\lean{Two.fixedIn}, \lean{Two.twoIn}), which with \lean{LemmaK} give \lean{ThmBab}
(\lean{thmBab_of_lemmaK}). For each bad prime $\ell\in B_\psi$, decomposing the character set
$X(e')\subseteq X(0)\cup X(1)$ (\lean{FinalLe.lean}) into odd-order characters $\mathcal O$
(at most $\mathrm{cond}(\psi)$ characters, each bounded by $B+M^{C_4}$) and remaining characters
(each bounded by $M^{2C_m}$), plus the finite-index loss $\le tM^6$, establishes hypothesis~(b) of
Lemma~\ref{lem:pointHV} (\lean{SumIn}, from \lean{CountII}, \lean{CountV}, \lean{countVs},
\lean{countIV}). Together with \lean{SumOut} (\lean{CountI}, \lean{CountZ}, \lean{CountII},
\lean{CountIII}), Lemma~\ref{lem:pointHV} proves \lean{PointHV}, which instantiates to \lean{ThmHV}
(\lean{ThmHV.lean}) with $B=\exp(x^{4\eps/5})$.

Fix $X\ge X_0$ and a window pair $(q,d)$ at $(\eps,X^{1/J})$, and partition $[X^{1/J},X]$ into
$\lfloor\log_2 J\rfloor$ interior-disjoint dyadic-exponent windows $[Y_r,Y_r^2]$ with
$Y_{r+1}=Y_r^2$. Combining \lean{ThmBab} and \lean{ThmHV} (\lean{jointWin_of_counts}) bounds the
number of exceptional scales where the plus-part (and $\ell=2$) sum exceeds
$(F\cdot t)^{C_1}+C_1\exp(x^{4\eps/5})$ by a constant independent of $X$. For $J\ge2$ sufficiently large, the pigeonhole principle yields a window $[Y,Y^2]$ avoiding all
exceptional scales, on which
\[
\log 2\sum_{1\le e'<t}\len\Sha_2(e')+\sum_{\ell\text{ odd}}\log\ell\sum_{1\le e'<t}\len\Sha_\ell^+(e')\le
(F\cdot t)^{C_1}+C_1\exp(x^{4\eps/5}).
\]
Since $\log F\ll x^{\eps/5}(\log x)^2$ and $t\le x^{401}$, both $(F\cdot t)^{C_1}$ and
$C_1\exp(x^{4\eps/5})$ are $o(\exp(x^\eps/2))$. Adding Steps~1 and~2 gives \eqref{eq:VVle}.
\end{proof}

% ======================================================================
\section{Proof of Theorem~C2}\label{sec:assembly}
% ======================================================================

\begin{proof}[Proof of Theorem~C2]
Fix a primitive Dirichlet character $\psi$ of order $o$, put $\eps_1=1/52$, and let
$0<\eps<\eps_1$. Let $e(x)=\lceil(2+8\eps)\log_2 x\rceil$. First apply Theorem~\ref{thm:VWin}
(using \ref{C4}, \ref{C5}, \ref{C7}) to obtain $J\ge2$ and $X_0>0$. For each sufficiently large
$X\ge\max(X_0, Y_0(\psi,\eps,J)^J)$, Proposition~\ref{prop:hostB} applied with this $J$ and
$Y_{\mathrm{base}}=X^{1/J}$ (using \ref{C1}) supplies a window pair $(q,d)$ at $(\eps,X^{1/J})$
whose auxiliary characters satisfy the zero-density box bounds on the full interval
$[X^{1/J},X]=[Y_{\mathrm{base}},Y_{\mathrm{base}}^J]$. Theorem~\ref{thm:VWin} then produces a
subwindow $[Y,Y^2]\subseteq[X^{1/J},X]$ on which the $\Sha$-window bound \eqref{eq:VVle} holds for
every $x\in[Y,Y^2]$. Letting $W$ be the union of these windows $[Y,Y^2]$ over all large $X$, we
obtain part~(i) of Theorem~C2 (since $[Y,Y^{3/2}]\subseteq[Y,Y^2]$ for $Y\ge1$).

Now fix $x\in W$ with $x\ge x_0(\eps)$. By Proposition~\ref{prop:ranks} and
Proposition~\ref{prop:hostB} (using \ref{C2}), we obtain a host at scale $x$ satisfying
\textup{(B0)--(B7)} with parameters $(C_E,K_0,\frac1{3o+1},4o)$, $G_K=G_{F_\psi}$,
$\Sigma_{\mathrm{rest}}(x)\le 4o\,x^{1/2+\eps}$, the domain conditions \textup{(S0)}, and all
windows of Proposition~7.9 of Part~I (with absolute constants $K,e_K$ coming from the parameter
bounds of Propositions~\ref{prop:ranks} and~\ref{prop:hostB}). Using the glued common lattice $R$
of Theorem~\ref{thm:gluedR} and the unitary theta-lift chain of \ref{C8}, we verify the ten slope
inequalities:
\begin{itemize}
\item \textup{(S1)} by Theorem~\ref{thm:S1} (using \ref{C3}, \ref{C6}, \ref{C7}, \ref{C8}, and
  \eqref{eq:VVle});
\item \textup{(S2)--(S5)} by Proposition~\ref{prop:S2S5};
\item \textup{(S6)} by Proposition~\ref{prop:S6};
\item \textup{(S7)} by Theorem~\ref{thm:S7};
\item \textup{(S8)--(S10)} by Proposition~\ref{prop:S8S10}.
\end{itemize}
The exponents $e_{\mathrm{Div}}$ and $e_L$ depend only on the absolute structural constants of the
chain and not on $\eps$, so setting $e=\max(e_{\mathrm{Div}}+e_L+3,3)$ completes the proof of
Theorem~C2. Combined with Theorems~A, B, and C1 of Part~I~\cite{dixit-part1}, this yields GRH for
all Dirichlet $L$-functions.
\end{proof}

% ======================================================================
\section{Literature-fidelity audit of inputs (C1)--(C8)}\label{sec:audit}
% ======================================================================

Every input \ref{C1}--\ref{C8} is proved internally inside the Lean~4 development using only the
three standard axioms (\texttt{propext}, \texttt{Quot.sound}, \texttt{Classical.choice}), so the
formal correctness of GRH does not depend on whether any internal statement matches its textbook
namesake at full strength. For a conventional mathematical reading, however, it is essential to
verify how each formal statement aligns with the published literature. Table~\ref{tab:litfidelity}
records the 1-to-1 correspondence and fidelity notes from our \lean{PropAudit} scan~\cite{dixit-repo}.

\begin{small}
\begin{longtable}{@{}p{11mm}p{45mm}p{46mm}p{44mm}@{}}
\caption{Literature-fidelity audit of classical inputs \ref{C1}--\ref{C7} and automorphic package \ref{C8}.}\label{tab:litfidelity}\\
\toprule
Input & Lean declaration(s) & Published reference & Fidelity / scope note \\
\midrule
\endfirsthead
\toprule
Input & Lean declaration(s) & Published reference & Fidelity / scope note \\
\midrule
\endhead
\ref{C1} & \lean{ZD.zero_density_box}, \lean{count_bad_core} & Iwaniec--Kowalski~\cite[Ch.~10]{iwaniec-kowalski} & Proved in Lean via mollifier/duality; matches classical $Q^{1/4}$ bad-modulus bound. \\
\ref{C2} & \lean{hostThickExists} & Behrend~\cite{behrend}, sphere-packing construction & Elementary combinatorial proof in Lean; gives density $\ge\exp(-C\sqrt{\log n\log L_{\mathrm{rel}}})/L_{\mathrm{rel}}$. \\
\ref{C3} & \lean{meanPeriod} & Chowla--Selberg~\cite{chowla-selberg} & Identifies $\log\Omega=\gE+\log(\pi/v)-L'(1,\chi_d)/L(1,\chi_d)$ for CM curves over $\Q(\sqrt{-d})$. \\
\ref{C4} & \lean{UpperTheory.V.FW.ferreroWashington} & Ferrero--Washington~\cite{ferrero-washington} & Proved in 62 Lean files; states bounded $\ell$-rank of $\mathrm{Cl}(K')/\mathrm{Cl}(K')^\ell$ ($\mu_\ell=0$). \\
\ref{C5} & \lean{PoitouTate.ShaDual}, \lean{LitCoh}, \lean{Odd.thmBcFree_of} & Neukirch--Schmidt--Wingberg~\cite[Ch.~VIII]{nsw}; Milne~\cite{milne} & Proved and applied via \lean{PropAudit}; covers minus-part orthogonality and Euler characteristic. \\
\ref{C6} & \lean{BlochKato.DimFormula}, \lean{RationalDuality}, \lean{IntegralDuality} & Bloch--Kato~\cite[Prop.~3.8, Cor.~3.8.4]{bloch-kato} & Proved via \lean{bkAL}/\lean{bkBL}; applied at crystalline places with $H^0=0$. \\
\ref{C7} & \lean{GrothendieckMessing}, \lean{GrothendieckDeformation}, \lean{CountI}--\lean{CountV} & Messing~\cite[Ch.~V, Thm.~1.6]{messing}; Grothendieck~\cite[Ch.~V~\S4]{grothendieck1974}; Diamond~\cite[\S2.4]{diamond} & Messing requires $p$ locally nilpotent (covering \lean{GrothendieckMessing}); \lean{GrothendieckDeformation} drops $p$-nilpotency, matching Grothendieck's 1974 Montr\'eal notes / Diamond (2025). \\
\ref{C8} & \lean{ShapeJumps}, \lean{BesselId}, \lean{ClassNumberFormula}$_2$, \lean{LinkIRR}, \lean{theoremS_link_concrete} & Family-specific unitary theta lifts, seesaw identities, class number formula & Certified in the live Lean closure; reduction steps unpacked in Sections~\ref{sec:glued}--\ref{sec:S8VWin}. \\
\bottomrule
\end{longtable}
\end{small}

% ======================================================================
\section*{Acknowledgements}
% ======================================================================

Every idea in the formal proof is due to Levent Alp\"oge~\cite{alpoge-remark}. This companion paper
was extracted directly from the formal Lean~4 archive (commit \texttt{40e56bab}) and builds on the
public verification work of Andrew Sutherland~\cite{sutherland}, Bartosz Naskr\k{e}cki, and the
contributors to \texttt{leanprover-community/riemann-hypothesis} (in particular \texttt{@Deicyde}
for pointing out the distinction between Messing's $p$-nilpotent theorem and Grothendieck's general
divided-power deformation statement).

\paragraph{Use of AI tools.} AI coding and research assistants were used to search the formal Lean
source and verify correspondences. All mathematical statements, deductions, and correspondences
were checked by the author, who takes full responsibility for the content.

\paragraph{License.} This paper is distributed under CC BY 4.0.

% ======================================================================
\begin{thebibliography}{99}
% ======================================================================
\bibitem{alpoge-remark} L.~Alp\"oge, note accompanying the Lean development, 20 November 2026,
  \url{https://alpo.ge/remark.pdf}.
\bibitem{behrend} F.~A.~Behrend, \emph{On sets of integers which contain no three terms in
  arithmetical progression}, Proc. Natl. Acad. Sci. USA \textbf{32} (1946), 331--332.
\bibitem{bloch-kato} S.~Bloch and K.~Kato, \emph{$L$-functions and Tamagawa numbers of motives},
  in: The Grothendieck Festschrift, Vol.~I, Progr. Math. \textbf{86}, Birkh\"auser, Boston, 1990,
  333--400.
\bibitem{chowla-selberg} S.~Chowla and A.~Selberg, \emph{On Epstein's zeta-function}, J. reine
  angew. Math. \textbf{227} (1967), 86--110.
\bibitem{diamond} F.~Diamond, \emph{Mod $p$ Hilbert modular forms and partial Hasse invariants},
  Forum Math. Sigma (2025), arXiv:2111.10160.
\bibitem{dixit-part1} I.~Dixit, \emph{The analytic reduction in Alp\"oge's formal proof of the
  Generalized Riemann Hypothesis}, preprint, 23 November 2026,
  \url{https://doi.org/10.5281/zenodo.24477213}.
\bibitem{dixit-repo} I.~Dixit, \emph{Reproducing and examining the Lean proof of the Riemann
  Hypothesis}, GitHub repository, version 1.2, 23 November 2026,
  \url{https://github.com/ishaan2006/rh-lean-reproduction}.
\bibitem{ferrero-washington} B.~Ferrero and L.~C.~Washington, \emph{The Iwasawa invariant $\mu_p$
  vanishes for abelian number fields}, Ann. of Math. (2) \textbf{109} (1979), 377--395.
\bibitem{grothendieck1974} A.~Grothendieck, \emph{Groupes de Barsotti--Tate et cristaux de
  Dieudonn\'e}, S\'eminaire de Math\'ematiques Sup\'erieures \textbf{45}, Presses de l'Universit\'e
  de Montr\'eal, 1974.
\bibitem{iwaniec-kowalski} H.~Iwaniec and E.~Kowalski, \emph{Analytic Number Theory}, AMS
  Colloquium Publications \textbf{53}, American Mathematical Society, Providence, RI, 2004.
\bibitem{messing} W.~Messing, \emph{The Crystals Associated to Barsotti--Tate Groups: With
  Applications to Abelian Schemes}, Lecture Notes in Mathematics \textbf{264}, Springer, 1972.
\bibitem{milne} J.~S.~Milne, \emph{Arithmetic Duality Theorems}, 2nd ed., BookSurge, Charleston,
  SC, 2006.
\bibitem{nsw} J.~Neukirch, A.~Schmidt and K.~Wingberg, \emph{Cohomology of Number Fields}, 2nd ed.,
  Grundlehren der mathematischen Wissenschaften \textbf{323}, Springer, Berlin, 2008.
\bibitem{sutherland} A.~V.~Sutherland, build and verification page for the Lean development,
  \url{https://math.mit.edu/~drew/grh/}, November 2026.
\end{thebibliography}

\appendix

% ======================================================================
\section{Correspondence with the formal development}\label{app:lean}
% ======================================================================

All declarations below belong to the live dependency closure of \lean{riemannHypothesis} in commit
\texttt{40e56bab} (recorded in \lean{results/used_decls.txt.gz} of~\cite{dixit-repo}). Paths are
relative to \texttt{src/}, with \texttt{Dev/AVT} abbreviating
\lean{Development/AbelianVarietyTheory}, \texttt{Dev/AA} abbreviating
\lean{Development/AutomorphicArithmetic}, \texttt{SWB} abbreviating
\lean{Formalism/SizeWindowBrackets}, and \texttt{SR} abbreviating
\lean{Formalism/SharpRestriction}.

\begin{small}
\begin{longtable}{@{}p{28mm}p{64mm}p{56mm}@{}}
\toprule
Here & Formal declaration & Location \\
\midrule
\endhead
Lemma~\ref{lem:evenbetween}, Prop.~\ref{prop:ranks} & \lean{exists_even_between}, \lean{exists_ranks_of_frame}, \lean{exists_ranks_rankVal_frame}, \lean{RankVal} & \lean{Dev/AVT/Proofs/GlobalInterpolation/ExistsRanksRankValFrame.lean}, \lean{PairRowIffQdRows.lean} \\
Prop.~\ref{prop:ranks}(iii) & \lean{m_pow_five_mul_f_mul_L_le}, \lean{size_window_at_five}, \lean{winL}, \lean{winA}, \lean{winB}, \lean{winM} & \lean{MPowFiveMulFMulLLe.lean}, \lean{QuadraticLogDerivCongr.lean} \\
Def.~\ref{def:winpair}, Prop.~\ref{prop:hostB} & \lean{WinPair}, \lean{PageN}, \lean{QH.pairEx}, \lean{count_bad_core}, \lean{ZD.zero_density_box}, \lean{hostThickExists} & \lean{Dev/AA/CharacterBoxes/ExistsPageN.lean} \\
Section~\ref{sec:glued}, \eqref{eq:sigchain}--\eqref{eq:rowranks} & \lean{IsChiField}, \lean{Ranks}, \lean{nuRow}, \lean{nuAbs} & \lean{ExistsPageN.lean}, \lean{Formalism/LocalRingCardinality/Params.lean} \\
Theorem~\ref{thm:gluedR} & \lean{prescribedRowLat}, \lean{exists_divStage_of_pk}, \lean{DV.RRDef}, \lean{GoodSideLattices.locMeet} & \lean{Dev/AVT/Proofs/IwasawaDecomposition/ExistsDivStageOfPk.lean:54, 76} \\
Lemmas~\ref{lem:heightchain}--\ref{lem:linkNgd}, Thm.~\ref{thm:S1} & \lean{height_chain_atR}, \lean{divClausesR_of}, \lean{errB_of_units}, \lean{wordFacts_of}, \lean{rank_le_col}, \lean{rank_le_row}, \lean{etaJ_le}, \lean{linkN_of_gd}, \lean{LinkDIV_at} & \lean{ConstantGeTwentyEight.lean:174}, \lean{RecBridge.lean:180}, \lean{JumpAndLengthOfAdmitted.lean:89, 115}, \lean{IsHeckeGaloisIff.lean:103}, \lean{LinkNOfGd.lean:77} \\
Lemmas~\ref{lem:lettersplace}--\ref{lem:C2off}, Prop.~\ref{prop:S2S5} & \lean{letters_place}, \lean{nuAbs_place_fine}, \lean{degree_fine}, \lean{C2off_of}, \lean{secondComplement_of_inputs}, \lean{C2atG_of} & \lean{FockDegreeGapBound.lean:204}, \lean{NuAbsPlaceFine.lean:101, 136, 177}, \lean{WeightDegreeSizes.lean:400, 460} \\
Prop.~\ref{prop:S6} & \lean{NprimeMuR}, \lean{LinkMu}, \lean{Places_at}, \lean{slopesN_at_of} & \lean{SR/GlobalTheoremInputs.lean:69}, \lean{SR/GlobalSchurBounds/TopDensityNGeSixteen.lean:1529} \\
Lemmas~\ref{lem:kWord}--\ref{lem:orbitsum}, Thm.~\ref{thm:S7} & \lean{LinkG_at}, \lean{key_minor_dvd}, \lean{kW_ord}, \lean{key3Ord}, \lean{caseD_arith}, \lean{caseE_arith}, \lean{St.starAll_of}, \lean{gain_bound_all_of_star} & \lean{KeyMinorDvd.lean:181, 207, 273}, \lean{PartnerCountOfCentral.lean:539, 587, 884}, \lean{NotSingularOfCentral.lean:777} \\
Lemmas~\ref{lem:diagnondeg}--\ref{lem:linkS}, Prop.~\ref{prop:S8S10} & \lean{unfolding_constant_pos_of_diagonal}, \lean{dall_pos}, \lean{mass_corollary}, \lean{proviso_e}, \lean{RankTiesF}, \lean{ClassNumberFormula}$_2$, \lean{VolSigIndep}, \lean{weight_numbers_bound}, \lean{epsilon_bound_large_x}, \lean{delp_of_allRows}, \lean{linkIRR_of_allRows}, \lean{locShare_le_Wg}, \lean{sizesRec_at}, \lean{closed_linkIRR_of_join_count}, \lean{hsNormSq_add_le}, \lean{diag_unproj}, \lean{linkS_concl}, \lean{size_thresholds}, \lean{eF_nonneg_and_le_exp_neg}, \lean{linkS_of_standing_hypotheses}, \lean{ClKppOff}--\lean{ClDiag}, \lean{matKF_of}, \lean{LinkS_at}, \lean{LinkIRR}, \lean{ChosenProviso}, \lean{theoremS_link_concrete} & \lean{OneAddInvPowLeExpOne.lean:56, 64}, \lean{MainSup.lean:135, 189}, \lean{ClassNumberClosed.lean:48}, \texttt{VolumeSignatureIndep\allowbreak AndClassNumberItems.lean:53}, \lean{LwRankTies.lean:54}, \lean{ConcNP.lean:59, 102, 276, 321, 459}, \lean{ConcNWeight.lean:284, 391}, \lean{LinkIRROfJoinEx.lean:61, 83}, \lean{LinkIrr.lean:60}, \texttt{StandingHypotheses\allowbreak HoldFamily.lean:75, 110, 171, 213, 233, 598}, \lean{ThmSOwn.lean:76}, \lean{RowsPos.lean:221, 261}, \lean{LinkSFinal.lean:57}, \lean{TopDensityNGeSixteen.lean:1343, 1357, 1411}, \lean{SR/GlobalTheoremInputs.lean:64}, \lean{LinkSPos.lean:151} \\
Lemmas~\ref{lem:minuspart}--\ref{lem:pointHV}, Thm.~\ref{thm:VWin} & \lean{VWin}, \lean{VWinAt}, \lean{VVle}, \lean{vwin_of_counts}, \lean{ThmBcFree}, \lean{point_arith}, \lean{thmHV}, \lean{thmBab_of_lemmaK}, \lean{UpperTheory.V.FW.ferreroWashington} & \lean{ExistsPageN.lean:44}, \lean{ThmBcOfPsiBound.lean:220, 339}, \lean{PointArith.lean:47, 78}, \lean{FinalLe.lean:199}, \lean{ThmHV.lean:49} \\
Section~\ref{sec:assembly} & \lean{ChainTied}, \lean{TieGlobal}, \lean{TieLocal}, \lean{TieArith}, \lean{supplyChi}, \lean{vwin_of_row} & \lean{SR/GlobalSchurBounds/TieGlobal.lean:97} \\
\bottomrule
\end{longtable}
\end{small}

\end{document}
